% ZXTune_LHUNA — reference manual
% Build: lualatex lhuna-bible-en (twice, for the table of contents)
% Formats tables are generated: see ../tools/formats.py --lang=en
\documentclass[11pt,a4paper,twoside,openright]{book}

\usepackage{fontspec}
\usepackage[english]{babel}
\usepackage[margin=1.5cm,bindingoffset=0.5cm]{geometry}
\usepackage{graphicx}
\usepackage{longtable}
\usepackage{booktabs}
\usepackage{array}
\usepackage{xcolor}
\usepackage[most]{tcolorbox}
\usepackage{fancyhdr}
\usepackage{titlesec}
\usepackage{enumitem}
\usepackage[hidelinks]{hyperref}
\usepackage{xspace}
% software name; the brand and the logo remain "LHUNA"
\newcommand{\soft}{ZXTune\_LHUNA\xspace}

\graphicspath{{../images/}}

\definecolor{lhunacyan}{RGB}{0,229,255}
\definecolor{lhunanight}{RGB}{11,14,20}
\definecolor{lhunagrey}{RGB}{110,118,130}

\setlist{itemsep=2pt,topsep=4pt}
\titleformat{\chapter}[display]{\normalfont\huge\bfseries}{\color{lhunacyan}\chaptertitlename\ \thechapter}{12pt}{\Huge}
\titleformat{\section}{\normalfont\Large\bfseries}{\color{lhunacyan}\thesection}{1em}{}

\pagestyle{fancy}
\fancyhf{}
\fancyhead[LE]{\small\leftmark}
\fancyhead[RO]{\small\rightmark}
\fancyfoot[C]{\small\thepage}
\renewcommand{\headrulewidth}{0.3pt}

% part to be written (search for "TO WRITE")
\newcommand{\arediger}[1]{\begin{tcolorbox}[colback=red!5,colframe=red!60!black,title=TO WRITE,breakable]#1\end{tcolorbox}}

% centered screenshot with caption, width relative to line
\newcommand{\capture}[3][0.45]{%
  \begin{figure}[htbp]\centering
  \includegraphics[width=#1\linewidth]{#2}
  \caption{#3}
  \end{figure}}

% two screenshots side by side
\newcommand{\capturedouble}[4]{%
  \begin{figure}[htbp]\centering
  \includegraphics[width=0.45\linewidth]{#1}\hfill
  \includegraphics[width=0.45\linewidth]{#2}
  \caption{#3 (left) and #4 (right)}
  \end{figure}}

% generated by tools/formats.py, do not edit

\newcommand{\lhunaModulesCount}{287}
\newcommand{\lhunaContainersCount}{97}
\newcommand{\lhunaFormatsTable}{%
\begin{longtable}{@{}l p{0.46\textwidth} l l@{}}
\toprule \textbf{Id} & \textbf{Format} & \textbf{Type} & \textbf{Family} \\ \midrule \endhead
AS0 & ASC Sound Master v0.x & Tracker module & AY-3-8910/YM2149 \\
ASC & ASC Sound Master v1.x & Tracker module & AY-3-8910/YM2149 \\
AYC & CPC AYC & Register stream & AY-3-8910/YM2149 \\
FTC & Fast Tracker v1.x & Tracker module & AY-3-8910/YM2149 \\
GTR & Global Tracker v1.x & Tracker module & AY-3-8910/YM2149 \\
PSC & Pro Sound Creator v1.xx & Tracker module & AY-3-8910/YM2149 \\
PSG & Programmable Sound Generator & Register stream & AY-3-8910/YM2149 \\
PSM & Pro Sound Maker & Tracker module & AY-3-8910/YM2149 \\
PT1 & Pro Tracker v1.x & Tracker module & AY-3-8910/YM2149 \\
PT2 & Pro Tracker v2.x & Tracker module & AY-3-8910/YM2149 \\
SQT & SQ-Tracker Compiled & Tracker module & AY-3-8910/YM2149 \\
ST1 & Sound Tracker v1.x & Tracker module & AY-3-8910/YM2149 \\
ST3 & Sound Tracker v3.x Compiled & Tracker module & AY-3-8910/YM2149 \\
STC & Sound Tracker v1.x Compiled & Tracker module & AY-3-8910/YM2149 \\
STP & Sound Tracker Pro Compiled & Tracker module & AY-3-8910/YM2149 \\
TXT & VortexTracker II & Tracker module & AY-3-8910/YM2149 \\
VTX & VTX (Vortex Project) & Register stream & AY-3-8910/YM2149 \\
YM & YM (ST-Sound Project) & Register stream & AY-3-8910/YM2149 \\
YM & YM (ST-Sound Project) Packed & Register stream & AY-3-8910/YM2149 \\
AY & AY/EMUL & Memory dump & AY-3-8910/YM2149, Beeper \\
PT3 & Pro Tracker v3.x & Tracker module & AY-3-8910/YM2149, TurboSound \\
RMT & Raster Music Tracker & Memory dump & Atari POKEY \\
SAP & Slight Atari Player Sound Format & Memory dump & Atari POKEY \\
DSF & Dreamcast Sound Format & Memory dump & Console: Dreamcast \\
GSF & GameBoy Advance Sound Format & Memory dump & Console: Game Boy Advance \\
USF & Ultra64 Sound Format & Memory dump & Console: Nintendo 64 \\
2SF & Nintendo DS Sound Format & Memory dump & Console: Nintendo DS \\
NCSF & Nitro Composer Sound Format & Memory dump & Console: Nintendo DS \\
PSF & Playstation Sound Format & Memory dump & Console: PlayStation \\
PSF2 & Playstation2 Sound Format & Memory dump & Console: PlayStation 2 \\
SSF & Sega Saturn Sound Format & Memory dump & Console: Sega Saturn \\
SNSF & Super NES Sound Format & Memory dump & Console: Super Nintendo \\
GBS & GameBoy Sound & Memory dump & Game Boy \\
2DX9 & Konami 2DX9 & Register stream & Game audio stream \\
2MSF & Banpresto RIFF wrapper & Register stream & Game audio stream \\
AAC & tri-Ace ASKA engine & Register stream & Game audio stream \\
AAX & CRI AAX & Register stream & Game audio stream \\
ACB & CRI Atom Cursheet Binary & Register stream & Game audio stream \\
ACM & InterPlay ACM & Register stream & Game audio stream \\
ADP & Nintendo ADP (raw) & Register stream & Game audio stream \\
ADP & Nex Entertainment NXAP & Register stream & Game audio stream \\
ADP & Xilam DERF & Register stream & Game audio stream \\
ADS & Sony AudioStream & Register stream & Game audio stream \\
ADS & Midway ADS & Register stream & Game audio stream \\
ADX & CRI ADX & Register stream & Game audio stream \\
AHX & CRI AHX & Register stream & Game audio stream \\
AIX & CRI AIX & Register stream & Game audio stream \\
AKB & Square Enix AKB & Register stream & Game audio stream \\
AKB & Square Enix AKB2 & Register stream & Game audio stream \\
AST & Nintendo AST & Register stream & Game audio stream \\
AUD & Westwood Studios AUD & Register stream & Game audio stream \\
AUS & Capcom AUS & Register stream & Game audio stream \\
AWB & CRI Atom Wave Bank & Register stream & Game audio stream \\
AWC & Rockstar AWC & Register stream & Game audio stream \\
BAF & Bizzare Creations bank file & Register stream & Game audio stream \\
BAF & Bizzare Creations band file (bad rips) & Register stream & Game audio stream \\
BANK & FSB5 FEV & Register stream & Game audio stream \\
BAO & Ubisoft BAO (atomic) & Register stream & Game audio stream \\
BGW & Square Enix BGW & Register stream & Game audio stream \\
BIK & RAD Game Tools Bink & Register stream & Game audio stream \\
BMP & Konami BMP & Register stream & Game audio stream \\
BNK & Sony BNK (LE) & Register stream & Game audio stream \\
BNK & Sony BNK (BE) & Register stream & Game audio stream \\
BNK & Audiokinetic Wwise Soundbank Container & Register stream & Game audio stream \\
BNK & Electronic Arts BNK & Register stream & Game audio stream \\
BWAV & Nintendo BWAV & Register stream & Game audio stream \\
CAF & tri-Crescendo CAF & Register stream & Game audio stream \\
CKD & Ubisoft CKD & Register stream & Game audio stream \\
CKS & Cricket Audio Stream & Register stream & Game audio stream \\
CPK & CRI CPK container & Register stream & Game audio stream \\
CSTM & Nintendo CSTM & Register stream & Game audio stream \\
CWAV & Nintendo CWAV & Register stream & Game audio stream \\
DSP & Nintendo ADPCM & Register stream & Game audio stream \\
DSP & Nintendo ADPCM (LE) & Register stream & Game audio stream \\
DSP & Koei Tecmo WiiBGM & Register stream & Game audio stream \\
DSP & Namco CSTR & Register stream & Game audio stream \\
DSP & Retro Studios RS03 & Register stream & Game audio stream \\
DSPW & Capcom DSPW & Register stream & Game audio stream \\
DVI & Capcom IDVI & Register stream & Game audio stream \\
FSB & FMOD Sample Bank (v1-v5) & Register stream & Game audio stream \\
FSTM & Nintendo FSTM & Register stream & Game audio stream \\
FWAV & Nintendo FWAV & Register stream & Game audio stream \\
GCA & Terminal Reality GCA & Register stream & Game audio stream \\
GENH & Generic Header & Register stream & Game audio stream \\
H4M & Hudson HVQM4 & Register stream & Game audio stream \\
HCA & CRI HCA & Register stream & Game audio stream \\
HPS & HAL Labs HALPST & Register stream & Game audio stream \\
HWAS & Vicarious Visions HWAS & Register stream & Game audio stream \\
IDSP & Namco IDSP (NUB container) & Register stream & Game audio stream \\
IDSP & Various IDSP & Register stream & Game audio stream \\
IFS & Konami Arcade Games Container & Register stream & Game audio stream \\
IKM & MiCROViSiON container & Register stream & Game audio stream \\
IMUS & LucasArts iMUSE & Register stream & Game audio stream \\
ISD & Inti Creates OGG Vorbis & Register stream & Game audio stream \\
ISWS & Sumo Digital iSWS & Register stream & Game audio stream \\
KCEY & Konami KCE Yokohama & Register stream & Game audio stream \\
KTSS & Koei Tecmo KTSS & Register stream & Game audio stream \\
KVS & Koei Tecmo KVS & Register stream & Game audio stream \\
LOGG & Encrypted OGG Vorbis & Register stream & Game audio stream \\
LOGG & Corrupted OGG Vorbis & Register stream & Game audio stream \\
LOPUS & Capcom OPUS container & Register stream & Game audio stream \\
LWAV & Ubisoft RIFF & Register stream & Game audio stream \\
LWAV & Ubisoft LyN & Register stream & Game audio stream \\
M4A & Apple QuickTime Container & Register stream & Game audio stream \\
MAB & Square Enix MAB & Register stream & Game audio stream \\
MCA & Capcom MCA & Register stream & Game audio stream \\
MIC & Sony MultiStream MIC & Register stream & Game audio stream \\
MIC & Koei Tecmo MIC & Register stream & Game audio stream \\
MSF & Sony MultiStream File & Register stream & Game audio stream \\
MSF & Konami MSF & Register stream & Game audio stream \\
MTX & Matric Software container & Register stream & Game audio stream \\
MUL & Crystal Dynamics MUL & Register stream & Game audio stream \\
MUL & Crystal Dynamics MUL (BE) & Register stream & Game audio stream \\
MUSC & Krome MUSC & Register stream & Game audio stream \\
MUSX & Eurocom MUSX & Register stream & Game audio stream \\
NAAC & Namco AAC & Register stream & Game audio stream \\
NPS & Namco NPSF & Register stream & Game audio stream \\
NSOPUS & Nihon Falcom FDK & Register stream & Game audio stream \\
NUB & Namco nuSound container & Register stream & Game audio stream \\
NUS3 & Namco nuSound container v3 & Register stream & Game audio stream \\
NWA & VisualArt's NWA & Register stream & Game audio stream \\
NXA & Entergram NXA & Register stream & Game audio stream \\
OMA & Sony OpenMG audio & Register stream & Game audio stream \\
OPUS & Nintendo OPUS & Register stream & Game audio stream \\
OPUS & Various OPUS containers & Register stream & Game audio stream \\
P3D & Radical P3D & Register stream & Game audio stream \\
PDT & Hudson Splitted Stream container & Register stream & Game audio stream \\
PDT & Hudson Stream Container & Register stream & Game audio stream \\
PONA & Policenauts BGM & Register stream & Game audio stream \\
RAK & Ubisoft RAKI & Register stream & Game audio stream \\
RAK & Ubisoft RAKI (modified) & Register stream & Game audio stream \\
RFRM & Retro Studios RFRM & Register stream & Game audio stream \\
RPGMVO & Encrypted OGG Vorbis & Register stream & Game audio stream \\
RSD & Radical RSD & Register stream & Game audio stream \\
RSD & RedSparc RSD & Register stream & Game audio stream \\
RSTM & Nintendo RSTM & Register stream & Game audio stream \\
RWAV & Nintendo RWAV & Register stream & Game audio stream \\
RWS & RenderWare Stream & Register stream & Game audio stream \\
RXW & Sony XWS (bad rips) & Register stream & Game audio stream \\
SAB & Square Enix SAB & Register stream & Game audio stream \\
SAD & Procyon Studio SADL & Register stream & Game audio stream \\
SCD & Square Enix SCD & Register stream & Game audio stream \\
SCHL & Electronic Arts SCHL & Register stream & Game audio stream \\
SD9 & Konami SD9 & Register stream & Game audio stream \\
SDF & Beyond Reality SDF & Register stream & Game audio stream \\
SEB & Game Arts SEB & Register stream & Game audio stream \\
SEG & Stormfront SEG & Register stream & Game audio stream \\
SMK & RAD Game Tools SMACKER & Register stream & Game audio stream \\
SMP & Infernal Engine SMP & Register stream & Game audio stream \\
SNR & Electronic Arts SNR & Register stream & Game audio stream \\
SPS & Electronic Arts SPS & Register stream & Game audio stream \\
SPSD & Naomi SPSD & Register stream & Game audio stream \\
STER & Alchemy STER & Register stream & Game audio stream \\
STM & Intelligent Systems STM & Register stream & Game audio stream \\
STR & Sega Stream Asset Builder STR & Register stream & Game audio stream \\
STR & 3DO STR & Register stream & Game audio stream \\
STRM & Nintendo STRM & Register stream & Game audio stream \\
SVAG & Konami SVAG & Register stream & Game audio stream \\
SVAG & SNK SVAG & Register stream & Game audio stream \\
SWAV & Nintendo SWAV & Register stream & Game audio stream \\
THP & Nintendo THP & Register stream & Game audio stream \\
VAG & Sony VAG & Register stream & Game audio stream \\
VGS & Princess Soft VGS & Register stream & Game audio stream \\
VPK & SCE America VPK & Register stream & Game audio stream \\
VXN & Gameloft VXN & Register stream & Game audio stream \\
WAVE & EngineBlack WAVE & Register stream & Game audio stream \\
WEM & Audiokinetic Wwise & Register stream & Game audio stream \\
WMSF & Banpresto MSF wrapper & Register stream & Game audio stream \\
XA & Sony XA & Register stream & Game audio stream \\
XA & Sony XA (raw) & Register stream & Game audio stream \\
XA & Maxis XA & Register stream & Game audio stream \\
XMA & Microsoft XMA & Register stream & Game audio stream \\
XNB & Microsoft XNA Game Studio XNB & Register stream & Game audio stream \\
XVAG & Sony XVAG & Register stream & Game audio stream \\
XWB & Microsoft XACT Wave Bank & Register stream & Game audio stream \\
XWC & Starbreeze XWC & Register stream & Game audio stream \\
XWM & Microsoft XWMA & Register stream & Game audio stream \\
XWS & Sony XWS & Register stream & Game audio stream \\
XWV & feelplus XWAV & Register stream & Game audio stream \\
YDSP & Yuke DSP & Register stream & Game audio stream \\
ZSD & ZSD container & Register stream & Game audio stream \\
ZSS & Z-Axis ZSND & Register stream & Game audio stream \\
HES & Home Entertainment System & Memory dump & HuC6280 \\
MIDI & Standard MIDI File & Tracker module & MIDI (Roland MT-32) \\
SID & Commodore64 SID/RSID/PSID & Memory dump & MOS6581/8580 \\
NSF & NES Sound Format & Memory dump & Ricoh 2A03/2A07 \\
NSFE & Extended Nintendo Sound Format & Memory dump & Ricoh 2A03/2A07 \\
COP & E-Tracker v1.x & Tracker module & SAA1099 \\
SPC & SNES SPC700 & Memory dump & SPC700/S-DSP \\
GYM & Genesis YM2612 (unpacked) & Register stream & Several chips \\
KSS & KSS Music Format & Memory dump & Several chips \\
KSSX & KSS Extended Music Format & Memory dump & Several chips \\
MTC & Multitrack Container & Multichip & Several chips \\
S98 & Sound 98 & Register stream & Several chips \\
VGM & Video Game Music & Register stream & Several chips \\
AAC & Advanced Audio Coding (raw ADTS) & Register stream & Standard audio \\
AC3 & Dolby AC-3 & Register stream & Standard audio \\
AIFC & Apple Compressed Audio Interchange File Format & Register stream & Standard audio \\
AIFF & Apple Audio Interchange File Format & Register stream & Standard audio \\
APE & Monkey's Audio & Register stream & Standard audio \\
ASF & Windows Media Audio & Register stream & Standard audio \\
CAF & Apple Core Audio Format File & Register stream & Standard audio \\
FLAC & Free Lossless Audio Codec & Register stream & Standard audio \\
MP3 & MPEG Audio Layer & Register stream & Standard audio \\
MPC & Musepack & Register stream & Standard audio \\
OGG & OGG Vorbis & Register stream & Standard audio \\
OPUS & OGG Opus & Register stream & Standard audio \\
TAK & Tom's lossless Audio Kompressor & Register stream & Standard audio \\
WAV & Waveform Audio File Format & Register stream & Standard audio \\
667 & Composer 667 & Tracker module & Tracker (samples) \\
669 & 669 Composer / UNIS 669 & Tracker module & Tracker (samples) \\
AHX & Abyss' Highest Experience & Tracker module & Tracker (samples) \\
AMF & ASYLUM Music Format & Tracker module & Tracker (samples) \\
AMF & DSMI / Digital Sound And Music Interface & Tracker module & Tracker (samples) \\
AMS & Extreme's Tracker & Tracker module & Tracker (samples) \\
AMS & Velvet Studio & Tracker module & Tracker (samples) \\
C67 & Composer 670 & Tracker module & Tracker (samples) \\
CBA & Chuck Biscuits / Black Artist & Tracker module & Tracker (samples) \\
CHI & Chip Tracker & Tracker module & Tracker (samples) \\
DBM & DigiBooster Pro & Tracker module & Tracker (samples) \\
DBM & Digi Booster & Tracker module & Tracker (samples) \\
DMF & X-Tracker & Tracker module & Tracker (samples) \\
DMM & Digital Music Maker v1.x & Tracker module & Tracker (samples) \\
DSM & Digital Sound Interface Kit RIFF & Tracker module & Tracker (samples) \\
DST & Digital Studio & Tracker module & Tracker (samples) \\
DSYM & Digital Symphony & Tracker module & Tracker (samples) \\
DTM & Digital Tracker / Digital Home Studio & Tracker module & Tracker (samples) \\
DTT & Desktop Tracker & Tracker module & Tracker (samples) \\
EMOD & Quadra Composer & Tracker module & Tracker (samples) \\
ET1 & Extreme Tracker v1.x & Tracker module & Tracker (samples) \\
ETX & EasyTrax & Tracker module & Tracker (samples) \\
FAR & Farandole Composer & Tracker module & Tracker (samples) \\
FC & Future Composer 1.0 - 1.4 & Tracker module & Tracker (samples) \\
FMT & Davey W Taylor's FM Tracker & Tracker module & Tracker (samples) \\
FNK & Funktracker & Tracker module & Tracker (samples) \\
FTM & Face The Music & Tracker module & Tracker (samples) \\
GDM & BWSB Soundsystem & Tracker module & Tracker (samples) \\
GMC & Game Music Creator & Tracker module & Tracker (samples) \\
GTK & Graoumf Tracker 1/2 & Tracker module & Tracker (samples) \\
HVL & Hively Tracker & Tracker module & Tracker (samples) \\
IMF & Imago Orpheus & Tracker module & Tracker (samples) \\
IMS & Images Music System & Tracker module & Tracker (samples) \\
IT & Impulse Tracker & Tracker module & Tracker (samples) \\
J2B & Galaxy Sound System & Tracker module & Tracker (samples) \\
KRIS & ChipTracker & Tracker module & Tracker (samples) \\
LIQ & Liquid Tracker & Tracker module & Tracker (samples) \\
LIQ & Liquid Tracker NO & Tracker module & Tracker (samples) \\
MDL & Digitrakker & Tracker module & Tracker (samples) \\
MED & MED 1.12 MED2 & Tracker module & Tracker (samples) \\
MED & MED 2.00 MED3 & Tracker module & Tracker (samples) \\
MED & MED 2.10 MED4 & Tracker module & Tracker (samples) \\
MED & OctaMED / MED Soundstudio & Tracker module & Tracker (samples) \\
MO3 & Un4seen MO3 & Tracker module & Tracker (samples) \\
MOD & ProTracker 3.6 & Tracker module & Tracker (samples) \\
MOD & Generic MOD-compatible & Tracker module & Tracker (samples) \\
MOD & MnemoTroN SoundTracker (MOD-compatible) & Tracker module & Tracker (samples) \\
MOD & Ice Tracker (MOD-compatible) & Tracker module & Tracker (samples) \\
MOD & Ultimate Soundtracker / etc (MOD Compatible) & Tracker module & Tracker (samples) \\
MT2 & Mad Tracker 2.xx & Tracker module & Tracker (samples) \\
MTM & MultiTracker & Tracker module & Tracker (samples) \\
OKT & Oktalyzer & Tracker module & Tracker (samples) \\
PDT & ProDigi Tracker v0.0x & Tracker module & Tracker (samples) \\
PLM & Disorder Tracker 2 & Tracker module & Tracker (samples) \\
PSM & Epic MegaGames MASI & Tracker module & Tracker (samples) \\
PSM & Epic MegaGames MASI (Old Version) & Tracker module & Tracker (samples) \\
PTM & PolyTracker & Tracker module & Tracker (samples) \\
PUMA & PumaTracker & Tracker module & Tracker (samples) \\
RTM & Real Tracker 2 & Tracker module & Tracker (samples) \\
S3M & ScreamTracker 3 & Tracker module & Tracker (samples) \\
SFX & SoundFX 1.x & Tracker module & Tracker (samples) \\
SFX & SoundFX 2.0 / MultiMedia Sound & Tracker module & Tracker (samples) \\
SQD & SQ Digital Tracker & Tracker module & Tracker (samples) \\
STIM & Slamtilt & Tracker module & Tracker (samples) \\
STM & ScreamTracker 2 & Tracker module & Tracker (samples) \\
STP & Soundtracker Pro II & Tracker module & Tracker (samples) \\
STR & Sample Tracker v2.x & Tracker module & Tracker (samples) \\
STX & STMIK 0.2 & Tracker module & Tracker (samples) \\
SYMMOD & Symphonie & Tracker module & Tracker (samples) \\
TCB & TCB Tracker & Tracker module & Tracker (samples) \\
ULT & UltraTracker & Tracker module & Tracker (samples) \\
V2M & Farbrausch V2 Synthesizer System & Tracker module & Tracker (samples) \\
XM & FastTracker II & Tracker module & Tracker (samples) \\
TF0 & TFM Music Maker v0.1-1.2 & Tracker module & TurboFM (YM2203) \\
TFC & TurboFM Compiled Dump & Register stream & TurboFM (YM2203) \\
TFD & TurboFM Dump & Register stream & TurboFM (YM2203) \\
TFE & TFM Music Maker v1.3+ & Tracker module & TurboFM (YM2203) \\
TS & TurboSound & Multichip & TurboSound \\
\bottomrule \end{longtable}}
\newcommand{\lhunaContainersTable}{%
\begin{longtable}{@{}l p{0.55\textwidth} l@{}}
\toprule \textbf{Id} & \textbf{Container} & \textbf{Type} \\ \midrule \endhead
7ZIP & 7zip & Archive \\
GZIP & GZip & Archive \\
HOBETA & Hobeta & Archive \\
HRIP & Hrip (Hrust RiP archiver) & Archive \\
LHA & LHA (LHArc) & Archive \\
RAR & RAR & Archive \\
UMX & UMX (Unreal Music eXperience) & Archive \\
ZDATA & Zdata & Archive \\
ZIP & ZIP & Archive \\
ZXZIP & ZXZip & Archive \\
CC3 & CodeCruncher v3.x & Compressor \\
CC4 & CompressorCode v.4 by ZYX & Compressor \\
CC4PLUS & CompressorCode v.4+ by ZYX & Compressor \\
CHARPRES & CharPres & Compressor \\
DSQ & DataSqueezer v4.x & Compressor \\
ESV & ESV Cruncher & Compressor \\
GAM & GamePacker & Compressor \\
GAMPLUS & GamePacker+ & Compressor \\
HRUM & Hrum v3.x & Compressor \\
HRUST1 & Hrust v1.x & Compressor \\
HRUST2 & Hrust v2.1 & Compressor \\
HRUST23 & Hrust v2.3 & Compressor \\
LZH1 & LZH Compressor v1.4 & Compressor \\
LZH2 & LZH Compressor v2.4 & Compressor \\
LZS & ASC LZS Pack & Compressor \\
MEGALZ & MegaLZ & Compressor \\
MSP & MicroSpace Packer v1.x & Compressor \\
MUSE & MUSE Compressor & Compressor \\
PACK2 & Pack v2.x & Compressor \\
PCD61 & Powerfull Code Decreaser v6.1 & Compressor \\
PCD61i & Powerfull Code Decreaser v6.1i & Compressor \\
PCD62 & Powerfull Code Decreaser v6.2 & Compressor \\
TLZ & Turbo-LZ v1.x & Compressor \\
TLZP & Turbo\#LZ v1.x (internal) & Compressor \\
TRUSH & Trush Compressor v3.x & Compressor \\
DSK & DSK Image & Disk image \\
FDI & FDI (Full Disk Image) & Disk image \\
SCL & SCL (SINCLAIR) & Disk image \\
TD0 & TD0 (TeleDisk Image) & Disk image \\
TRD & TRD (TR-DOS) & Disk image \\
AAC & Multitrack tri-Ace ASKA engine & Multitrack file \\
ACB & Multitrack CRI Atom Cursheet Binary & Multitrack file \\
AKB & Multitrack Square Enix AKB2 & Multitrack file \\
ASF & Multitrack Windows Media Audio & Multitrack file \\
AWB & Multitrack CRI Atom Wave Bank & Multitrack file \\
AWC & Multitrack Rockstar AWC & Multitrack file \\
AY & Multi-AY/EMUL & Multitrack file \\
BAF & Multitrack Bizzare Creations bank file & Multitrack file \\
BAF & Multitrack Bizzare Creations band file (bad rips) & Multitrack file \\
BIK & Multitrack RAD Game Tools Bink & Multitrack file \\
BNK & Multitrack Sony BNK (LE) & Multitrack file \\
BNK & Multitrack Sony BNK (BE) & Multitrack file \\
BNK & Multitrack Audiokinetic Wwise Soundbank Container & Multitrack file \\
BNK & Multitrack Electronic Arts BNK & Multitrack file \\
CPK & Multitrack CRI CPK container & Multitrack file \\
FSB & Multitrack FMOD Sample Bank (v1-v5) & Multitrack file \\
GBS & Multitrack GameBoy Sound & Multitrack file \\
H4M & Multitrack Hudson HVQM4 & Multitrack file \\
HES & Multitrack Home Entertainment System & Multitrack file \\
IFS & Multitrack Konami Arcade Games Container & Multitrack file \\
KSSX & Multitrack KSS Extended Music Format & Multitrack file \\
M4A & Multitrack Apple QuickTime Container & Multitrack file \\
MAB & Multitrack Square Enix MAB & Multitrack file \\
MUSX & Multitrack Eurocom MUSX & Multitrack file \\
NSF & Multitrack NES Sound Format & Multitrack file \\
NSFE & Multitrack Extended Nintendo Sound Format & Multitrack file \\
NUB & Multitrack Namco nuSound container & Multitrack file \\
NUS3 & Multitrack Namco nuSound container v3 & Multitrack file \\
PDT & Multitrack Hudson Stream Container & Multitrack file \\
RWS & Multitrack RenderWare Stream & Multitrack file \\
RXW & Multitrack Sony XWS (bad rips) & Multitrack file \\
SAB & Multitrack Square Enix SAB & Multitrack file \\
SAP & Multitrack Slight Atari Player Sound Format & Multitrack file \\
SCD & Multitrack Square Enix SCD & Multitrack file \\
SID & Multitrack Commodore64 SID/RSID/PSID & Multitrack file \\
VXN & Multitrack Gameloft VXN & Multitrack file \\
XA & Multitrack Sony XA & Multitrack file \\
XA & Multitrack Sony XA (raw) & Multitrack file \\
XVAG & Multitrack Sony XVAG & Multitrack file \\
XWB & Multitrack Microsoft XACT Wave Bank & Multitrack file \\
XWM & Multitrack Microsoft XWMA & Multitrack file \\
XWS & Multitrack Sony XWS & Multitrack file \\
ZSS & Multitrack Z-Axis ZSND & Multitrack file \\
COMPILEDASC0 & ASC Sound Master v0.x player & Player decompiler \\
COMPILEDASC1 & ASC Sound Master v1.x player & Player decompiler \\
COMPILEDASC2 & ASC Sound Master v2.x player & Player decompiler \\
COMPILEDPT24 & Pro Tracker v2.40 Phantom Family player & Player decompiler \\
COMPILEDPTU13 & Pro Tracker Utility v1.3 player & Player decompiler \\
COMPILEDST3 & Sound Tracker v3.x Compiled player & Player decompiler \\
COMPILEDSTP1 & Sound Tracker Pro v1.x player & Player decompiler \\
COMPILEDSTP2 & Sound Tracker Pro v2.x player & Player decompiler \\
RAW & Raw scaner & Raw data scanner \\
SNA128 & SNA 128k & Snapshot \\
Z80V145 & Z80 v1.45 & Snapshot \\
Z80V20 & Z80 v2.x & Snapshot \\
Z80V30 & Z80 v3.x & Snapshot \\
ZXSTATE & SZX (ZX-State) & Snapshot \\
\bottomrule \end{longtable}}


\begin{document}

\frontmatter
\begin{titlepage}
\pagecolor{lhunanight}\color{white}
\centering
\vspace*{3cm}
\includegraphics[width=0.6\textwidth]{lhuna_logo}\par
\vspace{1.5cm}
{\Huge\bfseries ZXTune\_LHUNA\par}\vspace{0.4cm}{\LARGE Reference manual\par}
\vspace{0.5cm}
{\Large Sound chip music player for Android\par}
\vfill
{\large Document version: \today\par}
\vspace{1cm}
\end{titlepage}
\nopagecolor\color{black}

\chapter*{License and trademark}
\addcontentsline{toc}{chapter}{License and trademark}

\section*{Under the LHUNA license}
The name \emph{LHUNA}, the logo, the visual identity and this manual are the property of their author and are
distributed under the LHUNA license. Any reproduction of the logo or of this manual requires the author's consent.

\section*{Under a free software license}
The \soft application is free software distributed under the \textbf{GNU General Public License, version 3}. This
follows from the components it contains:
\begin{itemize}
  \item \soft is derived from \emph{ZXTune}, distributed under the GNU Lesser General Public License version~3.
        LGPL~v3 code can be combined with GPL~v3, not with GPL~v2.
  \item Several bundled emulation engines are licensed under GPL~v3 (\emph{Highly Theoretical/Experimental}
        Saturn and PlayStation cores) or GPL~v2 ``or later'' (libsidplayfp, z80ex, ASAP), compatible with GPL~v3.
  \item The other components use permissive licenses (MIT, BSD, ISC, zlib, MPL~2.0\dots), all compatible.
\end{itemize}
Anyone receiving the APK is entitled to the corresponding source code, and may modify and redistribute it under the
same terms. The full list of components and licenses is given in appendix~\ref{chap:credits}.

\section*{Proprietary drivers}
The author's music drivers (\emph{NebulaWeave} family) are not part of the base APK. They are distributed in a
separate APK under a proprietary license.
\begin{itemize}
  \item Without that APK, \soft works exactly as described in this manual.
  \item With that APK installed, \soft can use the drivers it provides.
\end{itemize}
For both licenses to coexist, the drivers APK must remain a separate work under GPL~v3: it runs in its own process and
talks to \soft through a documented interface (Android service, content provider, data). Loading its native library
into the \soft process would make the whole a single work, which would then have to be entirely under GPL~v3.

\section*{Trademarks}
Names of consoles, computers, chips and games cited in this manual belong to their respective owners and are used
only to describe compatibility.

\tableofcontents
\chapter*{Before you start}
\addcontentsline{toc}{chapter}{Before you start}

\section*{How this software was made}
We developed \soft with the help of an artificial intelligence model, Claude by Anthropic, used as a programming
tool. We would rather say so clearly, once, here. The split was simple: we decided what the software should do, what
it should show and how to show it; we tested every build on our own hardware, with our own music, and turned down
whatever did not suit us. The tool wrote the code, the tests and the control renders from our instructions. No
build reaches you without being compiled, tested automatically, then tried on a real phone.

The work described in this manual runs from 2 October to 24 November 2026.

\section*{What you will find here}
This manual explains how to use \soft, but also what happens inside each sound chip and how to read it on screen.
If you only want to listen to your music, the first two parts are enough. If you compose, study the hardware or want
to understand why a tune sounds the way it does, the rest is for you.

A few conventions: registers are written in hexadecimal with the \$ prefix, notes in English notation (concert A is
A4, at 440~Hz) and levels in decibels relative to the chip's maximum. Screenshots are renders of the application on
a screen 1080 pixels wide.


\mainmatter
\part{Getting started}
\chapter{Overview}

\section{What it is}
\soft is an Android music player dedicated to music of computers, consoles and arcade machines from the 1980s to the
2000s. This music is not recorded audio: it is programs and data driving a sound chip. To play it faithfully, \soft
reproduces the behaviour of these chips in software, which is emulation.

\section{Origin}
\soft is derived from \emph{ZXTune}, a free player created for ZX Spectrum music and later extended to many
platforms. \soft reuses its playback core and adds:
\begin{itemize}
  \item emulation engines chosen for accuracy (ares for the Super Nintendo, Ootake for the PC Engine, reSIDfp for the
        Commodore~64, Nuked for Yamaha FM chips);
  \item real-time visualisation of each voice and of the internal state of chips, chip by chip;
  \item a redesigned interface: multiple playlists, favourites, play-next queue, live settings.
\end{itemize}

\section{Features}
\begin{itemize}
  \item Plays \lhunaModulesCount{} music formats and opens \lhunaContainersCount{} container formats (archives,
        disk images, packers), see part~\ref{part:compat}.
  \item Browses the phone storage and large online catalogues.
  \item Shows one oscilloscope per voice, the SID filter response, and detailed dashboards for SID, FM, SNES,
        PC Engine, NES, MT-32 and Sega arcade PCM chips.
  \item Works with Android Auto and the lock screen.
\end{itemize}

\section{Requirements}
\begin{itemize}
  \item Android~4.1 or later for the full build.
  \item Oscilloscope: OpenGL ES, or Vulkan where available (smoother).
  \item About 16~MB for the application, plus your music.
\end{itemize}

\section{History}
\subsection{ZXTune}
ZXTune comes from the ZX Spectrum community, whose music scene produced tens of thousands of tunes for the
AY-3-8910 chip with trackers such as Sound Tracker, Pro Tracker or ASC Sound Master. Each tracker stored music in
its own format, and players of the time read only one. ZXTune, created by Vitamin, set out to read all of them, to
find them even inside disk images, archives or packed programs, and to replay them faithfully.

The project then opened to other platforms by integrating the best free libraries: libsidplayfp for the
Commodore~64, Game Music Emu for consoles, libopenmpt and libxmp for tracker modules, vgmstream for modern game
audio, libvgm for VGM logs. It exists as a command-line tool, a desktop GUI and an Android application.

\subsection{ZXTune\_LHUNA}
\soft is a branch of the ZXTune Android application with two goals:
\begin{description}
  \item[Fidelity] use for each chip the most accurate engine available, with a lighter engine selectable for slow
        devices and comparison;
  \item[Readability] show what each chip does, voice by voice and frame by frame, with views built around the
        synthesis of each chip rather than a generic view.
\end{description}
Everything displayed is the real state of the emulated chip at the current moment: no averaged or invented values.
Text does not repeat the graphics; it gives what graphics do not show: register values, notes, levels in dB.

\chapter{Installation and first launch}

\section{Installing the APK}
\soft is distributed as an APK file.
\begin{enumerate}
  \item Copy the file to the phone, or download it directly.
  \item Open it. Android asks to allow installation from this source: accept for this time.
  \item An update installs over the previous version; playlists and settings are kept.
\end{enumerate}

\section{Permissions}
\begin{description}
  \item[Storage] to read music on the phone. On recent Android versions, \soft asks for access to a specific folder
        rather than to the whole storage.
  \item[Notifications] for the playback notification and the progress of folders being added to playlists.
\end{description}

\section{Tour}
The application has three screens, reached from the bottom bar:
\begin{enumerate}
  \item \textbf{Now playing}: current track, visualisation and controls.
  \item \textbf{Playlists}.
  \item \textbf{Browser}: your files and online catalogues.
\end{enumerate}

\capturedouble{now-playing-en}{playlist-en}{Now playing}{Playlists}
\capture{browser-sources-en}{Browser: sources}

\chapter{The interface}

\section{The Now playing screen}
\capture{now-playing-en}{Now playing screen, information folded}

At the top of the screen, the toolbar holds the fullscreen switch, the live settings, the choice between the
oscilloscope and the filter response (for the SID), the dashboard toggle and the main menu. The heart of the screen
is the visualisation: one oscilloscope per voice, then the chip dashboard, described in
chapter~\ref{chap:dashboards}.

Below the visualisation, the information block gives the track's title, author, chip, format and original driver.
We fold it by default to leave the height to the visualisation; the arrow at the bottom left unfolds it. Then comes
the position in the track, then the row of controls: information, shuffle, previous, play and pause, next, repeat,
the $\equiv$+ button and the track menu.

The $\equiv$+ button is for filing the current track: you can add it to your favourites or to any playlist,
including a new one. The track menu, on the right, lets you share it, make it a ringtone or show all its properties.

\capture{now-playing-info-en}{Now playing screen, information unfolded}

\section{The Playlists screen}
\capture{playlist-en}{Playlists screen}
At the top, search filters the playlist shown; the buttons next to it give statistics, export, clear and sort. Just
below, the selector shows the current playlist and its number of tracks; tap it to switch playlists or manage them
(chapter~\ref{chap:listes}). Each row shows the chip family, title, author and duration. The handle on the right
moves a track.

\section{The Browser screen}
\capturedouble{browser-sources-en}{browser-folder-en}{Sources}{folder contents}
The path is shown at the top; tap one of its parts to jump straight back to that level. A long press on a file or
folder starts multiple selection.

\section{Away from the Now playing screen}
When you are in the playlists or the browser, a mini-player above the bottom bar reminds you of the current track
and lets you skip or pause without changing screens.

\section{Landscape, tablet, fullscreen}
\capture[0.9]{landscape-now-playing-en}{Now playing screen in landscape}
In landscape, the visualisation takes the width and the dashboard moves beside it. On a tablet, all three screens
make use of the available space.

\capture[0.9]{tablet-playlist-en}{On a tablet}

Fullscreen shows the oscilloscope alone; in portrait, the dashboard stays below it. We avoid the camera cutout at the
top and the rounded corners at the bottom in portrait, and the side cutout in landscape. A short tap leaves
fullscreen, a long press opens the live settings.


\part{Listening}
\chapter{The browser}

\section{Your files}
The browser opens the phone's storage and SD cards. On Android~10 and later, access goes through the system folder
picker: you point once at your music folder and \soft remembers it.

\soft opens archives (ZIP, RAR, 7z, LHA\dots) and disk images (TR-DOS, SCL\dots) as folders: you do not need to
unpack them. It also finds music hidden in packed programs or machine snapshots, which is common in ZX Spectrum scene
collections.

A file holding several tracks (SID, NSF, AY, KSS\dots) plays directly when you tap it: its subsongs follow one
another, then playback continues with the next file in the folder.

\section{Online catalogues}
The browser also gives live access to large public collections. The \textbf{High Voltage SID Collection} gathers
nearly all Commodore~64 music, sorted by composer, and provides the duration database we use. \textbf{The Mod
Archive}, \textbf{Modland} and \textbf{Amiga Music Preservation} cover Amiga and PC tracker modules; Modland sorts
them by format, AMP documents the musicians. \textbf{ZXArt}, \textbf{ZXTunes} and \textbf{Aygor} archive the ZX
Spectrum scene, demo-party competitions included. \textbf{ASMA} is the HVSC of 8-bit Atari. \textbf{VGMRips} offers
VGM recordings of console and arcade games, whole game at a time. \textbf{Joshw} and \textbf{Khinsider} offer
console rips and soundtracks, \textbf{OverClocked ReMix} reinterpretations, \textbf{Aminet} and \textbf{Scene.org}
the historical archives of the Amiga and the demoscene. A few online radios complete the list.

These catalogues are independent services: their content and availability depend on the people who run them. Tracks
you have already played stay in the cache and play again without a connection.

\section{Searching and selecting}
The magnifier filters the folder shown. To act on several items, long-press the first then tap the others. The
``Add'' button asks which playlist to file them in; folders are scanned in full, which can take a while for a large
collection (progress shows in the notifications). The ``Play next'' button puts the selection in the queue, to hear
it right after the current track.

\chapter{Playlists}\label{chap:listes}

\section{Multiple playlists}
\soft handles any number of playlists, plus two special ones:
\begin{description}
  \item[Default playlist] created at first launch, it takes over your previous playlist.
  \item[Favourites] filled by the \textbf{$\equiv$+} button $\rightarrow$ ``Add to favorites''.
\end{description}
The displayed playlist is also the played one. Switching in the selector saves the playlist left and loads the new
one; the current track keeps playing and ``next'' follows the new playlist.

\capture{playlist-en-top}{Playlist selector}

\section{Managing playlists}
The selector menu offers: switch, new playlist, rename, duplicate, delete, export as M3U, import M3U. Deleting the
displayed playlist switches to another one.

\section{Adding tracks}
\begin{itemize}
  \item From Now playing: \textbf{$\equiv$+} $\rightarrow$ ``Add to playlist\dots'', then a playlist or ``New
        playlist\dots''.
  \item From the browser: selection $\rightarrow$ Add.
  \item From a playlist: selection $\rightarrow$ ``Add to playlist\dots'' copies tracks to another playlist.
\end{itemize}

\section{Play next queue}\label{sec:file}
Queued tracks play after the current one, in the order they were added. Then playback resumes where it was in the
playlist or folder. Choosing another track yourself keeps the queue but drops that resume point.

\section{M3U import and export}
Export writes the displayed playlist as extended M3U (title, author, duration). Locations are complete, including
inside archives and online catalogues, so an exported playlist is re-imported identically. Import also accepts M3U
from other players (plain paths) and creates a new playlist named after the file. Former XSPF playlists remain
available in the browser.

\section{Sort, move, statistics}
Sorting is by title, author or duration. Each track's handle moves it. Statistics give track count, distinct files
and total duration of the playlist or the selection.

\chapter{Playback}

\section{Shuffle and repeat}
Shuffle picks the next track at random in the playlist or folder. Repeat plays the current track endlessly; without
it, the track stops at the end of its known duration.

\section{Duration, loops and fades}
Much game music loops forever: the file does not say when to stop. When a duration is known, from the format itself
or from a database (the HVSC \emph{Songlengths} database for SID), we use it. Otherwise we apply the default
duration, which you set in the settings, along with the loop limit and the fades. When a track has an introduction
before its loop, the progress bar marks where the looped part starts.

\section{Outside the application}
You can control playback from the notification, the lock screen, a Bluetooth headset or Android Auto. Android Auto
shows your playlists and a simplified visualisation.

\section{Sharing}
The track menu shares the file, sends a link to the track, makes a ringtone from an excerpt and shows every known
property of the track.


\part{Compatibility}\label{part:compat}
\chapter{Synthesis concepts}\label{chap:synthese}
Dashboards use the vocabulary of the chips. This chapter explains each term by what it does to the sound.

\section{Oscillator and waveform}
An oscillator produces a repeated pattern; its repetition rate is the pitch of the note, its shape is the timbre.
\begin{description}
  \item[Square, pulse] two alternating levels. The \emph{duty cycle} is the share of time spent at the high level: at
        50~\% the sound is hollow (odd harmonics only), at 12.5~\% it is nasal.
  \item[Triangle] linear rise and fall; few harmonics, a soft flute-like tone.
  \item[Sawtooth] linear rise then sharp drop; all harmonics, a bright brass or string tone.
  \item[Sine] pure tone, no harmonics.
  \item[Noise] pseudo-random sequence produced by a feedback shift register (LFSR); its clock sets the colour, from
        low rumble to high hiss. 8-bit chip percussion is made of noise.
\end{description}

\section{Frequency divider}
A chip does not compute a frequency: it divides its clock. For the AY, $f = F/(16N)$; for the PC Engine,
$f = F/(32N)$. The larger $N$, the lower the note. In practice, two neighbouring values of $N$ give widely spaced
notes in the treble, hence the poor tuning of high notes. The SID works the other way: a 24-bit phase accumulator
advances by $N$ each cycle, $f = N \cdot F / 2^{24}$, which gives it the same precision over the whole range.

\section{Envelope}
The envelope changes the volume of a note over time.
\begin{description}
  \item[ADSR] four phases: \emph{attack} (rise from silence), \emph{decay} (fall to the sustain level),
        \emph{sustain} (level held while the note is on), \emph{release} (fade after the note is released). A, D and
        R are durations, S is a level. The dashboard shows the current phase, the level and durations in
        milliseconds.
  \item[GAIN] (Super Nintendo) the program drives the volume directly or picks a linear or exponential ramp instead
        of the four phases.
  \item[AY hardware envelope] one envelope shared by the three voices, with a repeated pattern (sawtooth, triangle);
        at audible rates it becomes a waveform of its own.
\end{description}

\section{Filter}
A filter removes part of the spectrum.
\begin{description}
  \item[Low-pass] keeps the bass, dulls the treble; \emph{high-pass} the opposite; \emph{band-pass} keeps a band;
        \emph{notch} removes a band.
  \item[Cutoff frequency] the limit between what passes and what is attenuated.
  \item[Resonance] boost around the cutoff; pushed high, the filter whistles.
\end{description}
The SID is the only chip of this list with an adjustable analogue filter; the response curve shown by \soft gives,
for each frequency, the boost or attenuation applied, and the position of each voice.

\section{Sync and ring modulation}
\begin{description}
  \item[Sync] the oscillator is reset at each period of the source oscillator. Its waveform is cut at the source rate:
        the perceived pitch is the source's, the timbre changes with the voice's own frequency (``hard sync'' of
        aggressive bass and leads).
  \item[Ring] (ring modulation) the voice's triangle is multiplied by the sign of the source oscillator. The result
        contains the sum and difference of both frequencies: metallic tones, bells.
\end{description}
On the SID the source is fixed: the previous voice (V1 from V3, V2 from V1, V3 from V2). The Mod column shows it:
``Sync V3'', ``Sync+Ring V1''.

\section{FM synthesis}
Frequency modulation synthesis varies the frequency of a sine oscillator, the \emph{carrier}, at the rate of another,
the \emph{modulator}. The larger the modulator amplitude, the richer the harmonics.
\begin{description}
  \item[Operator] a sine oscillator with its own envelope and level (TL, \emph{total level}, attenuation in 0.75~dB
        steps).
  \item[Multiplier] ratio between the operator frequency and the note frequency; an integer ratio gives a harmonic
        tone, a detune (DT) a beating.
  \item[Algorithm] how operators are connected: in series (rich tone), in parallel (organ), mixed. Operators feeding
        the output are carriers, their level sets the volume. Those modulating others set the timbre.
  \item[Feedback] the first operator modulates itself: from sine towards sawtooth, then noise.
  \item[LFO] slow oscillator (a few hertz) adding vibrato (PM) or tremolo (AM).
\end{description}

\section{Sample playback}
\begin{description}
  \item[Sample] recorded sound, played at a rate that sets the pitch: at double rate, one octave higher.
  \item[Loop] after the end, playback restarts from a loop point to sustain a note.
  \item[BRR, ADPCM] sample compression (Super Nintendo: 9 bytes for 16 samples).
  \item[Interpolation] computing values between two samples; the Super Nintendo uses a Gaussian interpolation that
        softens the sound.
  \item[Wavetable] (PC Engine, Game Boy) a very short single-period sample written to memory by the program and
        looped: the table is the instrument.
  \item[DDA] (PC Engine) the processor writes the output value directly, to play long samples.
\end{description}

\section{Volume, panning, effects}
\begin{description}
  \item[Attenuation in dB] chips set volume in logarithmic steps; $-6$~dB halves the amplitude.
  \item[Panning] left/right balance, by separate volumes (PC Engine, Super Nintendo) or by switches (Mega Drive).
  \item[Echo and FIR] (Super Nintendo) sound is fed back after a delay through an 8-tap filter.
  \item[DSP] programmable processor of the Saturn and Sega arcade boards computing reverb and effects.
  \item[Samples through volume] (SID) changing the master volume very fast produces a 4-bit sample.
\end{description}

\chapter{Emulated platforms and chips}

\begin{longtable}{@{}p{0.22\textwidth} p{0.27\textwidth} p{0.43\textwidth}@{}}
\toprule \textbf{Platform} & \textbf{Chip} & \textbf{Synthesis} \\ \midrule \endhead
ZX Spectrum 128, Amstrad CPC, Atari ST, MSX & AY-3-8910 / YM2149 & 3 square voices, noise, hardware envelope \\
ZX Spectrum (TurboSound) & 2 $\times$ AY & 6 voices \\
ZX Spectrum (TurboFM) & 2 $\times$ YM2203 & 4-operator FM + AY \\
SAM Coupé & SAA1099 & 6 square voices, noise, envelopes \\
Commodore 64 & MOS 6581 / 8580 (SID), up to 3 chips & 3 oscillators, analogue filter, ring modulation, sync \\
NES / Famicom & Ricoh 2A03 / 2A07 & 2 pulses, triangle, noise, DPCM \\
Super Nintendo & SPC700 + S-DSP & 8 BRR sample voices, ADSR, FIR echo \\
Game Boy & LR35902 (APU) & 2 pulses, programmable wave, noise \\
PC Engine / TurboGrafx-16 & HuC6280 & 6 wavetable voices 32 $\times$ 5 bits, noise, DDA, LFO \\
Atari 8-bit & POKEY & 4 voices, noise polynomials, filters \\
Mega Drive / Genesis & YM2612 + SN76489 & 6-voice 4-operator FM, DAC, PSG \\
PC (AdLib, Sound Blaster) & YM3526 / YM3812 / YMF262 & 2- or 4-operator FM, rhythm mode \\
Sharp X68000, arcade & YM2151 & 8-voice 4-operator FM \\
Sega Model 1 and 2 arcade & YMW258 (MultiPCM) & 28 sample voices \\
Sega Saturn, Model 2A/3 arcade & YMF292 (SCSP) & 32 PCM/FM voices, effects DSP \\
Roland MT-32 & LA (Munt) & linear arithmetic synthesis \\
\bottomrule
\end{longtable}

\section{AY-3-8910 and YM2149 (ZX Spectrum, Amstrad CPC, Atari ST, MSX)}
The AY-3-8910 programmable sound generator was designed by General Instrument in the late 1970s. Yamaha produced a
compatible variant, the YM2149, used in the Atari ST. It is found in the ZX Spectrum 128, Amstrad CPC, MSX, Oric,
many arcade boards and expansion cards.

Voice frequency: $f = F_{clock} / (16 \times N)$, where $N$ is the 12-bit period. Volume is 4 bits (logarithmic, about
3~dB per step on the AY, finer steps on the YM2149). The envelope has a 16-bit period and eight shapes. The clock
depends on the machine: 1.7734~MHz for the ZX Spectrum, 1.75~MHz for the Pentagon, 2~MHz for the Atari ST, 1~MHz for
the Amstrad CPC. Composers built chords from fast arpeggios and used the envelope at audible rates as a fourth
waveform.

\soft reads the formats of the scene trackers (Pro Tracker, Sound Tracker, ASC, SQ-Tracker, Vortex Tracker\dots),
register logs (PSG, YM, VTX) and AY files containing the original player code, run on an emulated Z80.

\section{TurboSound and TurboFM}
Russian ZX Spectrum clones received expansion cards: TurboSound doubles the AY (six voices), TurboFM adds two Yamaha
YM2203 chips, each combining three FM voices and an AY. \soft then shows two columns of voices and the FM operators.

\section{SAA1099 (SAM Coupé)}
The Philips SAA1099 is used in the SAM Coupé, the British successor of the Spectrum, and in Creative Music System
cards. It produces six square channels, two noise generators and two envelopes, with per-channel stereo. E-Tracker is
the main format.

\section{SID (Commodore 64)}
The SID (Sound Interface Device) was designed by Bob Yannes at MOS Technology for the Commodore~64, released in 1982.
Designed as a real synthesizer, it was far ahead of other chips of its time and gave rise to a generation of
composers: Rob Hubbard, Martin Galway, Ben Daglish, Jeroen Tel, Chris Hülsbeck and many others.

Voice frequency: $f = N \times F_{clock} / 2^{24}$, with a 16-bit $N$; the clock is about 0.985~MHz in PAL and
1.023~MHz in NTSC. The pulse width is 12 bits. Sync resets the oscillator at the rate of the previous voice; ring
modulation multiplies the triangle by the previous voice's oscillator (V1 from V3, V2 from V1, V3 from V2). Register
\$D418 sets the 4-bit master volume and the filter mode.

\subsection{6581 and 8580}
The original 6581 has a filter that varies a lot from chip to chip and a current leak producing a click on each volume
change. The 8580, its replacement from 1986, has a more regular filter but sounds different; it strongly attenuates
samples played through the volume. \soft emulates both and offers filter curves measured on real chips.

\subsection{Samples on the SID}
The SID has no sample player, but composers invented two techniques:
\begin{description}
  \item[Volume register] writing the master volume (\$D418) thousands of times per second produces a 4-bit sample;
        it adds to the output, outside the three voices.
  \item[Test bit] freezing and restarting an oscillator lets a value be written into it; the sample then goes through
        a voice and its filter.
\end{description}
The samples lane of the oscilloscope shows the first technique separately.

\subsection{Several SIDs}
Some recent tunes use two or three chips (expansion cartridges, emulators). \soft places them left, right and centre,
or mixes them in mono depending on the setting.

\section{NES and Famicom (Ricoh 2A03)}
The NES processor, the Ricoh 2A03 (2A07 in Europe), contains its sound unit: two pulse channels with selectable duty
cycle (12.5~\%, 25~\%, 50~\%, 75~\%), a triangle, a noise channel and a DPCM channel playing 1-bit samples. Koji Kondo,
Hirokazu Tanaka and Tim Follin wrote classic soundtracks for it. \soft reads NSF and NSFe files, which contain the
game's player code.

\section{Super Nintendo (SPC700 and S-DSP)}
Super Nintendo sound is handled by a subsystem designed by Sony: an SPC700 processor and a signal processor, the
S-DSP, sharing 64~KB of memory. The S-DSP plays eight voices of BRR-compressed samples, with ADSR or GAIN envelope,
stereo volume and a FIR-filtered echo that gives the console's music its depth. It does not synthesize: it plays back
short instrument recordings, so sound quality depends on the samples chosen and on very limited memory.

An SPC file is a snapshot of the sound subsystem memory and registers; an SNSF file also holds the main processor
program, for games driving sound from it.

\section{Game Boy}
The Game Boy has two pulse channels (one with frequency sweep), a programmable wave channel of 32 4-bit steps and a
noise channel. GBS files contain the player code.

\section{PC Engine (HuC6280)}
NEC and Hudson Soft's PC Engine, called TurboGrafx-16 in the United States, integrates sound into its HuC6280
processor. Each of its six channels plays a 32-step 5-bit wavetable written by the program: this is wavetable
synthesis, where the waveform itself is the instrument. Channels 5 and 6 can produce noise; all can switch to DDA
mode, where the processor writes samples directly; channel 2 can modulate the frequency of channel 1 (LFO).

Registers are at \$0800: \$800 selects the channel, \$801 sets left and right master volume, \$802 and \$803 the
12-bit period, \$804 enable, DDA mode and 5-bit volume, \$805 left and right volume, \$806 wave data, \$807 noise,
\$808 and \$809 the LFO. Voice frequency: $f = 3\,579\,545 / (32 \times N)$.

\section{Atari 8-bit (POKEY)}
POKEY, designed by Doug Neubauer at Atari, is used in the Atari 400, 800, XL and XE and in many arcade machines. Its
four channels combine frequency dividers and noise polynomials, giving a very distinctive timbre; two channels can be
paired for better pitch precision. SAP files are played by the ASAP library.

\section{Mega Drive (YM2612 and SN76489)}
Sega's Mega Drive pairs a Yamaha YM2612 FM chip (six 4-operator voices, one able to play 8-bit samples) with the
SN76489 generator inherited from the Master System (three square voices and noise). Yuzo Koshiro, Masato Nakamura and
Motoi Sakuraba used all of its resources. Music is played from VGM logs.

\section{AdLib and Sound Blaster (OPL, OPL2, OPL3)}
On PC, AdLib then Sound Blaster cards use the Yamaha YM3812 (OPL2) then YMF262 (OPL3) FM chips. The OPL2 offers nine
2-operator voices, or six voices and five percussion sounds in rhythm mode; the OPL3 doubles the voice count, adds
4-operator voices and stereo. Since rhythm mode is rarely used, \soft shows the nine melodic voices in a 3~$\times$~3
grid by default.

\section{YM2151 (OPM)}
The OPM offers eight 4-operator voices and a rich LFO. It is found in the Sharp X68000, Yamaha synthesizers of the same
generation and many arcade boards.

\section{Sega Model 1, 2, 2A and 3 arcade boards}
Sega's 3D arcade boards use sample-based chips:
\begin{description}
  \item[Model 1 and Model 2] 68000 processor, YM3438 and two YMW258 (MultiPCM) with 28 sample voices each;
  \item[Model 2A, 2B, 2C and Model 3] SCSP (YMF292), the Saturn chip: 32 PCM voices that can FM-modulate each other,
        envelopes, LFO and a programmable effects DSP for reverb.
\end{description}

\section{Roland MT-32}
Roland's MT-32 MIDI module, widely used by late-1980s PC games, combines samples and synthesis in its LA (linear
arithmetic) synthesis. \soft emulates it with the Munt library, which requires images of the module's ROMs.

\chapter{Music formats}
\soft recognises \lhunaModulesCount{} music formats. The table below is generated from the application itself, so it
is exact for the version described.

\begin{description}
  \item[Tracker module] score and instruments, replayed by a player faithful to the original one;
  \item[Register stream] sequence of values written to the chip registers, logged from the machine;
  \item[Memory dump] original player program run on an emulated processor (Z80, 6502\dots);
  \item[Multichip] several chips played together.
\end{description}
``Game audio stream'' refers to audio formats of modern games (CRI, Nintendo, Sony, Wwise\dots): compressed recordings
read by the vgmstream library.

{\small\lhunaFormatsTable}

\section{Containers}
\soft also opens \lhunaContainersCount{} container formats: archives, packers, disk images, machine snapshots and
player decompilers.

{\small\lhunaContainersTable}

\chapter{Recorded audio formats}
Besides emulated music, \soft plays common audio formats:
\begin{itemize}
  \item WAV: 8 to 32-bit PCM, 32-bit floating point, extended header (WAVE\_FORMAT\_EXTENSIBLE);
  \item FLAC, Monkey's Audio, Musepack, TAK (lossless);
  \item MP3, Ogg Vorbis, Opus, AAC, AC-3;
  \item and game audio formats supported by vgmstream (previous chapter).
\end{itemize}
For these formats, Now playing shows codec, sample rate, bit depth and bitrate.
\capture{now-playing-info-flac-en}{Information of a FLAC file}


\part{Emulation engines}
\chapter{Emulation engines}

For several chips there are several emulators, more or less accurate and more or less demanding. We picked the most
faithful one as the default engine, and kept a lighter alternative where one exists. If your phone struggles (the
``render'' line of the statistics drops below 100~\%), change engine in the settings.

\section{Commodore 64: reSIDfp and reSID}
Both engines come from the libsidplayfp library, which also emulates the Commodore~64 processor and the circuits the
original players need to run. reSID is Dag Lem's historical model, built from measurements of real chips. reSIDfp,
developed later by Antti Lankila then Leandro Nini, models the analogue filter in floating point and reproduces
combined waveforms from samples of real chips; it is our default engine.

\section{Super Nintendo: ares and snes\_spc}
The ares engine, from the work of Near and their team on bsnes and higan, reproduces the SPC700 and the S-DSP cycle
by cycle. It renders stereo exactly as the console does, phase inversions included: some tracks therefore sound wider
than in other players, and that is expected. snes\_spc, written by Shay Green (blargg), is lighter and remains
available. The same setting applies to SPC and SNSF.

\section{Yamaha FM chips}
For each FM family you choose the core:
\begin{longtable}{@{}l l p{0.5\textwidth}@{}}
\toprule \textbf{Family} & \textbf{Chips} & \textbf{Cores} \\ \midrule
OPN2 & YM2612, YM3438 & Nuked (default), ymfm, Genesis Plus GX \\
OPM & YM2151 & Nuked (default), ymfm, MAME \\
OPN & YM2203, YM2608, YM2610 & ymfm (default), MAME \\
OPL & YM3526, YM3812, YMF262 & Nuked (default), AdLibEmu, MAME \\
OPLL & YM2413 & Nuked (default), EMU2413, MAME \\
\bottomrule
\end{longtable}
Alexey Khokholov's Nuked cores were written from a study of the chips' silicon; they are the most exact and the most
expensive. ymfm, by Aaron Giles, is the emulation MAME uses. Eke-Eke's Genesis Plus GX and MAME's older cores are
fast.

\section{PC Engine: Ootake}
We emulate the HuC6280 with the core of Ootake, Kitao Nakamura's PC Engine emulator, known for its fidelity to the
hardware. It is also from this core that we read each voice's wavetable for the PC Engine view.

\section{Saturn and Sega arcade: SCSP}
The SCSP is emulated with its effects processor on. We turned it on by default because the music of these machines
was mixed with its reverb: without it, it sounds dry and empty.

\section{What accuracy costs}
On a recent phone every engine runs well above real time. The most demanding are Nuked OPM and Nuked OPN2, then ares;
the lightest are Genesis Plus GX, the MAME cores and snes\_spc.

\chapter{Emulation settings: options and state of the art}\label{chap:emu-settings}

Every setting in this chapter changes the sound. For each one we give the possible choices, the default (in bold in
the tables), what you will hear when you change it, and what reference players and emulators do. When our default
differs from theirs, we explain why.

We invented almost none of these values. Filter presets come from published measurements of real chips and real
boards, and we name their origin each time. The raw list of every setting, with its internal key, is in
chapter~\ref{chap:all-settings}.

\section{General settings}
\begin{reglages}
\regl{Default duration}{30~s, 1, 2, \textbf{3}, 5, 10, 15, 30~min}{Used only when neither the file nor a database
gives the duration. Game music players usually choose between two and three minutes; three minutes let at least one
full loop of most tracks play.}
\regl{Silence detection}{\textbf{none}, 3, 4, 5, 10, 15~s}{Game Music Emu and the foobar2000 plugins built on it
skip to the next track after a silence. We leave it off: many tracks have deliberate pauses, and a SID playing
samples can seem silent for a few seconds.}
\regl{Loop limit}{\textbf{none}, 1 to 5}{For formats that know their loop point (VGM, some modules). Without a limit,
the known or default duration applies.}
\regl{Fade-in, fade-out}{\textbf{none}, 1, 2, 3, 5, 10~s}{A fade-out of a few seconds softens the cut of looping
tracks; VGMPlay and most VGM players offer one. We do not force it, so as never to alter a track that has a real
ending.}
\regl{Gain}{50 to 200~\%, \textbf{100~\%}}{Applied after emulation. Above 100~\%, tracks that are already loud clip:
turn up the phone volume instead.}
\end{reglages}

\section{SID (Commodore 64)}
\capture{settings-sid}{SID settings in the live panel}

The SID is the chip where the state of the art has gone furthest, and where settings matter most: two 6581s from the
same factory do not have the same filter. The software reference is the reSID family: Dag Lem's reSID, used by VICE,
and reSIDfp, developed later in libsidplayfp (sidplayfp on the command line) and ported to Java by JSIDPlay2.

\begin{reglages}
\regl{Engine}{reSID, \textbf{reSIDfp}}{sidplayfp uses reSIDfp by default, VICE offers reSID. reSIDfp models the 6581
filter from its transistors and reproduces combined waveforms better; reSID is lighter.}
\regl{Default model}{\textbf{6581}, 8580}{Used when the file does not say. The vast majority of HVSC music was written
on the 6581; it is also libsidplayfp's default.}
\regl{Force model}{\textbf{no}, yes}{To hear a track on the other chip. Useful for comparing, rarely faithful.}
\regl{C64 model}{\textbf{PAL}, NTSC, old NTSC, Drean (PAL-N), PAL-M}{The clock changes pitch and tempo. Most European
music is PAL; recent files state their model.}
\regl{Force C64 model}{\textbf{no}, yes}{As for the chip model.}
\regl{Interpolation}{\textbf{none}, performance, quality}{Maps to reSID's sampling methods: fast sampling,
interpolation, full resampling. Full resampling, which VICE also offers, removes treble aliasing but costs much more.
On a phone the difference is rarely audible.}
\regl{Filter emulation}{no, \textbf{yes}}{Without the filter, filtered voices sound raw: wrong for almost all C64
music. The setting is mostly there to understand what the filter brings.}
\regl{6581 preset}{\textbf{custom}, 31 presets}{See below.}
\regl{8580 preset}{\textbf{custom}, 5 presets}{See below.}
\regl{reSID: 6581 bias}{$-1000$ to $+1000$~mV, \textbf{0}}{Shifts the whole cutoff curve. VICE sets this bias to
$+500$~mV by default (\texttt{SidResidFilterBias}); we keep reSID's original value and offer VICE's as a preset.}
\regl{reSIDfp: 6581 curve}{0 (light) to 100 (dark), \textbf{50}}{Same scale and default as sidplayfp's
\texttt{FilterCurve6581} (0.5).}
\regl{reSIDfp: 6581 range}{0 (dark) to 100 (light), \textbf{50}}{Same default as sidplayfp.}
\regl{reSIDfp: 8580 curve}{0 (light) to 100 (dark), \textbf{50}}{Same default as sidplayfp.}
\regl{Combined waveforms}{\textbf{average}, weak, strong}{The strength of waveforms obtained by enabling several at
once varies from chip to chip. sidplayfp offers the same three levels (\texttt{CombinedWaveforms}), with the same
default.}
\regl{6581 pulse + saw}{\textbf{normal}, strong}{Some chips give a much louder pulse + sawtooth combination. Try it
when a sound seems missing in a track you know.}
\regl{8580 digi boost}{\textbf{no}, yes}{The 8580 makes samples played through the volume almost inaudible.
Musicians added a resistor to the board; libsidplayfp offers the same correction (\texttt{digiBoost}).}
\regl{2SID / 3SID output}{\textbf{stereo}, mono}{Stereo: one chip left, one right, the third centre, as most players
do. Mono: all chips centred, as on a real machine.}
\regl{Samples lane (\$D418)}{\textbf{no}, yes}{Specific to \soft. Separates samples played through the volume into
their own lane; the three voices then show only their oscillators. Applied to the next track.}
\end{reglages}

\subsection*{Filter presets}
The 31 6581 presets and the 5 8580 presets come from JSIDPlay2's default configuration. They fall into three
families.
\begin{description}
  \item[Lightest to Darkest] seven reSID curves, from lightest to darkest. JSIDPlay2 defines them by a curve
        position; we convert it to a reSID bias.
  \item[VICE default] VICE's $+500$~mV bias, to get that emulator's sound.
  \item[Measured chips] 23 reSIDfp curves named after the person who measured the chip and its date code: Alankila,
        Lord Nightmare, ZrX, Trurl, Nata, Grue, KruLLo, Enigma. ``6581R3 4885'' means an R3 revision made in week 48 of
        1985.
\end{description}
For the 8580: three reSID curves (Light, Average, Dark) and two chips measured by Trurl.

Our advice: leave the custom setting at its defaults, which are sidplayfp's. If a track sounds muffled, try Light or a
Trurl chip; if it whistles, Dark. There is no ``right'' 6581: the composer heard their own.

\section{Super Nintendo (SPC, SNSF)}
\capture{settings-spc}{SNES settings}
\begin{reglages}
\regl{Engine}{snes\_spc, \textbf{ares}}{ares carries on the work of bsnes and higan, the accuracy reference; Mesen is
the other emulator at that level. blargg's snes\_spc, used by Game Music Emu and many players, is lighter but does
not render the phase-inverted stereo some games use to widen the sound. Applied to the next track.}
\end{reglages}
We offer no interpolation other than Gaussian. Snes9x, bsnes and Mesen offer cubic or sinc interpolation, which makes
the sound brighter; pleasant, but not the console's sound, whose S-DSP applies exactly that Gaussian filter.

\section{Yamaha FM chips}
\capturedouble{settings-ym2612-psg}{settings-ym3812}{OPN2 settings}{OPL settings}
The core of an FM chip is digital: emulated well, it gives the same result as the chip. What changes from one
machine to another is what comes after, the DAC and the output circuit. That is why each family has a platform
setting, which sets the chip, low-pass filter, DC blocking and stereo from published measurements. If you change one
of those settings, the platform switches to ``custom''.

\subsection*{YM2612 and YM3438 (OPN2, Mega Drive)}
\begin{reglages}
\regl{Core}{\textbf{Nuked}, ymfm, Genesis Plus GX}{Nuked-OPN2 is written from the silicon and is the reference;
Genesis Plus GX also includes it as an option. ymfm is MAME's core. Genesis Plus GX's own core is the fastest.}
\regl{Platform}{\textbf{Mega Drive model 1 (VA0--VA6)}, model 1 VA7, model 2, Genesis 3, bare YM3438, bare YMF276,
custom}{Early Mega Drives have a separate YM2612; from model 1 VA7, the chip is built into the main ASIC as a YM3438.
Genesis Plus GX offers the same YM2612 / YM3438 choice.}
\regl{Chip}{\textbf{YM2612}, YM3438, YMF276}{Set by the platform.}
\regl{DAC ladder effect}{no, \textbf{yes} (model 1)}{Crossover distortion of the separate YM2612, which boosts quiet
sounds. Nuked-OPN2 reproduces it in YM2612 mode. The YM3438 lacks it.}
\regl{Low-pass filter}{0 to 24\,000~Hz, \textbf{3390~Hz}}{Model 1 filters its output with a first-order RC circuit.
BlastEm uses 3390~Hz; Genesis Plus GX's one-pole filter, set to 0.6, cuts around 3.7~kHz. Model 2 filters more: a
second order around 2840~Hz, measured by consolemods.org.}
\regl{Slope}{\textbf{6}, 12~dB/oct}{First order for model 1, second order for model 2.}
\regl{DC blocking}{0 to 200~Hz, \textbf{5~Hz}}{The output capacitors remove the DAC's DC voltage; this is essential
with the ladder effect, which offsets every voice even when silent. The real value is undocumented, so we put it
below the audible band.}
\regl{Stereo width}{0 to 100~\%, \textbf{100~\%}}{FM voices are only left, right or centre; reducing the width
helps on headphones with very separated tracks.}
\end{reglages}

\subsection*{YM2151 (OPM)}
\begin{reglages}
\regl{Core}{Nuked, \textbf{ymfm}, old MAME}{Nuked-OPM is the most exact but also the most expensive of all our cores:
we leave it as an option and use ymfm, MAME's current core.}
\regl{Platform}{\textbf{YM3012} (stereo: X68000, arcade), YM3014 (mono), custom}{The floating-point DAC is emulated
by the cores. Boards have no explicit filter on the FM part (noted in jotego's arcade work), so we add none.}
\regl{Filter, slope, DC, stereo}{\textbf{0~Hz}, \textbf{12~dB/oct}, \textbf{0~Hz}, \textbf{100~\%}}{No filtering by
default, for lack of measurements.}
\end{reglages}

\subsection*{YM2203, YM2608, YM2610 (OPN)}
\begin{reglages}
\regl{Core}{\textbf{ymfm}, old MAME}{There is no Nuked core for these chips. ymfm is MAME's core.}
\regl{Platform}{\textbf{PC-88 / PC-98}, custom}{The output stages of PC-88 and PC-98 boards have not been measured;
we do not filter.}
\regl{Clock (TurboFM)}{1.75, \textbf{3.5}, 3.58, 4~MHz}{For ZX Spectrum TurboFM music, the board runs at 3.5~MHz.}
\end{reglages}

\subsection*{YM3526, YM3812, YMF262 (OPL)}
\begin{reglages}
\regl{Core}{\textbf{Nuked}, AdLibEmu, MAME}{Nuked-OPL3 is the reference today, used in particular by modern DOSBox
versions. Ken Silverman's AdLibEmu and MAME's core are fast.}
\regl{Platform}{AdLib (mono), Sound Blaster 1.x/2.0, SB Pro 1/2, \textbf{SB16}, custom}{Values taken from DOSBox
Staging, which tuned them on recordings of real cards: SB 1.x and 2.0, first-order low-pass at 12~kHz; SB Pro,
first order at 8~kHz; SB16, no filter. The AdLib card has not been measured.}
\end{reglages}

\subsection*{YM2413 (OPLL)}
\begin{reglages}
\regl{Core}{\textbf{Nuked}, EMU2413, MAME}{Mitsutaka Okazaki's EMU2413 is the most widespread core in MSX and Master
System players and emulators; Nuked-OPLL is more exact.}
\end{reglages}

\section{AY-3-8910, YM2149 and the Z80 processor}
\capture{settings-ay}{AY settings}
\begin{reglages}
\regl{Clock}{ZX Spectrum (1.7734~MHz), \textbf{Pentagon (1.75~MHz)}, Taganrog (3.5~MHz), Atari ST (2~MHz), Amstrad CPC
(1~MHz)}{Most ZX scene music archived on ZXArt was written on the Pentagon, the most widespread clone in Russia, hence
our default. Files that state their clock take precedence.}
\regl{Type}{\textbf{AY-3-8910}, YM2149}{The YM2149 has a 32-step envelope instead of 16 and a slightly different
volume table. Choose it for the Atari ST.}
\regl{Interpolation}{none, performance, \textbf{quality}}{Smooths the conversion to 48~kHz. Without it the sound is
rawer, which some prefer on the Spectrum.}
\regl{Layout}{\textbf{ABC}, ACB, BAC, BCA, CAB, CBA, mono}{The original ZX Spectrum 128 is mono; stereo boards imposed
ABC or ACB depending on the region. Fuse and Unreal Speccy also offer ABC, ACB and mono.}
\regl{Mixer}{A: 100~\% left; B: 60~\% both sides; C: 100~\% right}{The three-channel mixer sets each voice. The centre
at 60~\% keeps a balanced level in mono.}
\regl{Duty cycle}{channels: \textbf{none} to all; 1 to 99~\%, \textbf{50~\%}}{A ZXTune option that matches no
hardware: it replaces the square with a pulse. Use it out of curiosity.}
\regl{Z80 clock}{\textbf{3.5~MHz}, 3.5469~MHz (128K), 4~MHz (CPC), 7~MHz}{For AY files that contain the player
code.}
\regl{Interrupt duration}{1 to 100~T, \textbf{24~T}}{How long the interrupt stays active each frame. Change it only if
a player misses interrupts.}
\end{reglages}

\section{SAA1099}
\begin{reglages}
\regl{Clock}{\textbf{SAM Coupé (8~MHz)}, Creative Music System (7.16~MHz)}{The two machines that use it.}
\regl{Interpolation}{none, \textbf{performance}, quality}{As for the AY.}
\end{reglages}

\section{Roland MT-32}
\begin{reglages}
\regl{ROM folder}{to choose}{Munt, also used by DOSBox and ScummVM, needs images of the module's control and sound
memories. We recognise them by their contents, whole or in halves.}
\regl{Preferred model}{\textbf{automatic}, MT-32, CM-32L / LAPC-I}{The CM-32L has extra sounds later games use; the
first games were written for the MT-32.}
\regl{Oscilloscope lanes}{\textbf{9 parts}, 32 partials}{Parts match the instruments; partials, the LA synthesis
building blocks. The second view will mostly interest those studying the module.}
\end{reglages}

\section{What we do not offer, and why}
\begin{itemize}
  \item No equaliser or ``surround'': they transform the machine's sound.
  \item No non-Gaussian interpolation for the Super Nintendo, for the same reason.
  \item No invented filter for boards that have not been measured (OPM, OPN, AdLib): we would rather add nothing than
        add a guess. As soon as a measurement is published, we will include it.
\end{itemize}


\part{Visualisation}
\chapter{The oscilloscope}

\section{One lane per voice}
The oscilloscope draws the shape of each voice's sound. Each chip takes a column, each voice a row, with its name at
the top of the cell. You see who plays what: the bass that holds, the arpeggio that shimmers, the snare that flashes
past.

To keep shapes still, we lock each lane's window to the period of its note: this is triggering, as on a lab
oscilloscope. We then reduce the signal to two points per pixel column, the minimum and the maximum, so no peak is
lost even when the lane is narrow.

\section{Special cases}
With the SID samples lane on, a fourth row per chip shows the samples played through the volume. For the AdLib, we lay
out the nine melodic voices 3 by 3 and hide the five rhythm-mode lanes, which few tracks use; turn on ``AdLib rhythm
lanes'' if you need them.

\section{SID filter response}
The visualisation button replaces the top half of the oscilloscope with the filter curve of each SID, with each
voice's position. Previous states stay visible for a moment as they fade, which makes fast cutoff sweeps readable.
\capturedouble{ux4-lowpass-waves}{ux4-sweep}{Resonant low-pass filter}{cutoff sweep}
If you want more room for the voices, reduce the filter response share in the oscilloscope settings.
\capture[0.7]{ux4-lowpass-share35}{Filter response reduced to 35~\% of the height}

\section{Settings}
\capture{prefs-en-oscilloscope}{Oscilloscope settings}
\begin{reglages}
\regl{Waveform window}{5 to 100~ms, \textbf{40~ms}}{Length of sound shown per voice. 40~ms is the value of corrscope
(\texttt{render\_ms}), the tool behind most chiptune oscilloscope videos.}
\regl{Visual gain}{\textbf{automatic per voice}, $\times$1 to $\times$8}{Affects only the drawing. In automatic mode
each lane is scaled to its own level, so a quiet voice stays readable.}
\regl{Maximum periods shown}{off, 4, \textbf{8}, 16, 32}{Zooms in on high notes, which otherwise shrink to a thick
line. Specific to \soft.}
\regl{SID gauges window}{2 to 40~ms, \textbf{10~ms}}{Length summarised by the SID oscillator gauges.}
\regl{Oscilloscope height}{60 to 160~\%, \textbf{100~\%}}{Relative to a 16:9 shape of the width; the dashboard takes
the rest.}
\regl{Filter response height}{20 to 80~\%, \textbf{55~\%}}{Share of the SID filter view taken by the curve.}
\regl{Waveforms under the filter}{no, \textbf{yes}}{Otherwise the curve alone.}
\regl{AdLib rhythm lanes}{\textbf{no}, yes}{See above. Applied to the next track.}
\regl{Frame rate}{\textbf{screen rate}, 120, 90, 60, 30}{A lower limit saves battery.}
\regl{Graphics engine}{\textbf{automatic}, Vulkan, OpenGL ES}{Automatic picks Vulkan when available.}
\regl{Custom Vulkan driver}{\textbf{no}, package to choose}{On Adreno GPUs, a Mesa Turnip driver can replace the
vendor's, as recent Android emulators do. The system driver takes over if yours fails to load.}
\end{reglages}

\section{Triggering}
Locking a waveform so it stays still is harder than it looks: a note changing timbre, an arpeggio, noise have no clean
edge. Lab oscilloscopes look for a threshold crossing; that is not enough for music. The state of the art is
corrscope's correlation trigger, by nyanpasu64: it looks, in the incoming sound, for the position most similar to the
previous frames, favouring rising edges. We ported it as is, with its default values, so that people who know
corrscope get the same picture and can reuse their settings.
\begin{reglages}
\regl{Preset}{\textbf{corrscope default (40~ms)}, evolving notes (60~ms), chiptune square, basses, NES triangle,
power saving, custom}{Each preset sets the settings below.}
\regl{Width}{10 to 200~ms, \textbf{40~ms}}{corrscope's \texttt{trigger\_ms}. Must cover two periods of the lowest
notes: raise it for basses.}
\regl{Edge strength}{0 to 300~\%, \textbf{100~\%}}{Favours steep rising edges: new notes are caught better, held
notes move more.}
\regl{Slope width}{5 to 100~\% of the period, \textbf{25~\%}}{Width of the edge looked for.}
\regl{Buffer strength}{0 to 300~\%, \textbf{100~\%}}{Favours similarity with previous frames.}
\regl{Responsiveness}{0 to 100~\%, \textbf{50~\%}}{How fast the history follows the sound. corrscope expresses it per
frame at 60~fps; we convert it for other screen rates.}
\regl{Falloff}{10 to 200~\% of the period, \textbf{50~\%}}{Weight of the history's edges.}
\regl{Reset below}{0 to 100~\%, \textbf{30~\%}}{Forgets the history when the new note no longer resembles it.}
\regl{DC removal}{0 to 100~\%, \textbf{0}}{For lanes whose centre drifts.}
\regl{Sign amplification}{0 to 300~\%, \textbf{0}}{Creates sharp zero crossings for soft waves such as the triangle.}
\regl{Search range}{10 to 100~\% of the width, \textbf{50~\%}}{How far the shape may jump from one frame to the
next.}
\regl{Falling edges}{\textbf{no}, yes}{For sounds that start with a fall, like some sampled trumpets.}
\regl{Pitch tracking}{no, \textbf{yes}}{Rescales the history when the note changes.}
\regl{Zero-crossing snap}{\textbf{0} samples}{Slides to the nearest rising zero crossing.}
\regl{Full searches}{\textbf{every frame}, 60, 30 per second}{Fewer searches, less computation.}
\regl{Subsampling}{\textbf{automatic}, none, 2, 4}{Computes the correlation on fewer points.}
\end{reglages}

\chapter{Chip dashboards}\label{chap:dashboards}

Below the oscilloscope, the dashboard details the chip state. Each chip has its own view, built around what is
specific to its synthesis. State snapshots are taken about every millisecond of emulated sound; each frame shows the
latest one, without smoothing.

\section{SID (Commodore 64)}
\capture[0.8]{sidcells-now-playing}{SID dashboard}
For each voice: waveform and duty cycle, envelope (level, phase, ADSR in milliseconds), frequency (note and hertz).
The \textbf{Mod} column shows sync and ring modulation with the source voice: ``Sync V3'', ``Ring V2'', and
``Sync+Ring V1'' when both are on. On the SID a voice is always modulated by the previous one (V1 from V3, V2 from V1,
V3 from V2). The last row shows volume, resonance and filter.
\capturedouble{sidcells-dual}{sidcells-fullscreen}{Two SID chips}{fullscreen}

\section{Super Nintendo (S-DSP)}
\capturedouble{snes}{snes-inspector}{View of the 8 voices}{voice inspector}
For each voice: sample, frequency, drawn ADSR or GAIN envelope, left and right volumes, echo, noise, pitch
modulation. At the bottom: master volumes, echo and its FIR filter. Tapping a voice opens its inspector.

\section{PC Engine (HuC6280)}
\capture[0.8]{pce-dashboard}{PC Engine view}
For each of the 6 voices: the 32-step wavetable in use, a level bar and two panning bars, then note, frequency and
level. The waveform shape is named (sine, triangle, rising saw, falling saw, square, pulse with its duty cycle) and
each distinct table gets a number, which changes when the instrument changes. The bottom row gives the registers as
written (\$802 to \$807); the global row, \$801, \$808 and \$809. The HuC6280 has no reverse playback: a falling saw is
a table written backwards by the driver, hence ``Saw $\downarrow$ (reversed)''.

\section{FM chips}
\capturedouble{fm-megadrive-cards}{fm-megadrive-inspector}{Mega Drive: voice cards}{voice inspector}
\capturedouble{fmview-algorithms}{fm-adlib-cards}{FM algorithms}{AdLib}
Each voice shows its algorithm (carriers and modulators), levels and envelopes of its operators. Tapping a voice opens
the inspector: drawn envelopes, multipliers, detune, LFO.

\section{NES, MT-32 and PCM chips}
\capturedouble{nes}{mt32parts}{NES}{Roland MT-32: parts}
\capturedouble{pcmchips-model1-portrait}{pcmchips-scsp-portrait}{YMW258 (Model 1/2)}{SCSP (Saturn, Model 2A/3)}

\chapter{Reading the values}

A few conventions hold for every view; once you know them, you can read any dashboard.

\section{What you see is the exact state}
We show no averages and no smoothing. An averaged value is a value the chip never had, and it hides exactly what is
interesting: an arpeggio changing note every frame, an envelope restarting. If a value flickers, the music makes it
flicker.

\section{Notes and frequencies}
Notes are written in English: C, D, E, F, G, A, B, with \# for sharps. Octave 4 holds the 440~Hz A (A4), middle C
is C4. When the frequency falls between two notes we show the nearest; the frequency in hertz gives you the gap. A
slight detune between two voices playing the same note is often deliberate: it thickens the sound.

\section{Levels}
Levels are relative to the chip's maximum: 0~dB is full volume, $-6$~dB half the amplitude, $-\infty$ silence. The
bars follow the same scale.

\section{Registers}
Registers are hexadecimal, prefixed with \$, as in period documentation: \$D418 on the SID, \$804 on the PC Engine.
We show them as the driver wrote them, not as we interpret them; if you look them up in a datasheet you will find the
same values.

\chapter{Live settings panel}
A long press on the oscilloscope, or the settings button, opens a panel over the visualisation. It shows only the
settings relevant to the current track: settings of the chip being played, of the oscilloscope and general settings.
Changes are heard immediately unless stated otherwise.
\capturedouble{settings-sid}{settings-general}{SID page}{General page}


\part{Reference}
\chapter{All settings}\label{chap:all-settings}
This chapter is generated from the application: it lists every setting in the order of the settings screen, with its
choices, default value and internal key (useful in a bug report).
\capturedouble{prefs-en-top}{prefs-en-3}{Settings, top}{continued}
% generated by tools/settings.py, do not edit
\section{Settings}
\subsection*{Other}
\begin{description}
\item[Default duration]  \emph{Choices:} 30 seconds, 1 minute, 2 minutes, 3 minutes, 5 minutes, 10 minutes, 15 minutes, 30 minutes. \emph{Default:} 3 minutes. {\footnotesize\texttt{zxtune.core.plugins.default\_duration}}
\item[Silence detection]  \emph{Choices:} None, 3 seconds, 4 seconds, 5 seconds, 10 seconds, 15 seconds. \emph{Default:} None. {\footnotesize\texttt{zxtune.sound.silencelimit}}
\item[Loops limit]  \emph{Choices:} None, 1 loop, 2 loops, 3 loops, 4 loops, 5 loops. \emph{Default:} None. {\footnotesize\texttt{zxtune.sound.looplimit}}
\item[Fade-in]  \emph{Choices:} None, 1 second, 2 seconds, 3 seconds, 5 seconds, 10 seconds. \emph{Default:} None. {\footnotesize\texttt{zxtune.sound.fadein}}
\item[Fade-out]  \emph{Choices:} None, 1 second, 2 seconds, 3 seconds, 5 seconds, 10 seconds. \emph{Default:} None. {\footnotesize\texttt{zxtune.sound.fadeout}}
\item[Gain] (range: 50\% to 200\%) \emph{Default:} 100\%. {\footnotesize\texttt{zxtune.sound.gain}}
\end{description}
\subsection*{AY–3–8910/YM2149F}
\begin{description}
\item[Clockrate]  \emph{Choices:} ZX Spectrum (1.7734 MHz), Pentagon (1.75 MHz), Taganrog (3.5 MHz), Atari ST (2 MHz), Amstrad CPC (1 MHz). \emph{Default:} Pentagon (1.75 MHz). {\footnotesize\texttt{zxtune.core.aym.clockrate}}
\item[Type]  \emph{Choices:} AY–3–8910, YM2149F. \emph{Default:} AY–3–8910. {\footnotesize\texttt{zxtune.core.aym.type}}
\item[Interpolation]  \emph{Choices:} None, Performance, Quality (high CPU usage). \emph{Default:} Quality (high CPU usage). {\footnotesize\texttt{zxtune.core.aym.interpolation}}
\item[Channels layout]  \emph{Choices:} ABC, ACB, BAC, BCA, CAB, CBA, Mono. \emph{Default:} ABC. {\footnotesize\texttt{zxtune.core.aym.layout}}
\item[Duty cycle channels]  \emph{Choices:} None, A, B, C, A, B, A, C, B, C, All. \emph{Default:} None. {\footnotesize\texttt{zxtune.core.aym.duty\_cycle\_mask}}
\item[Duty cycle] Tone pulse width for selected channels, 50 - square wave (range: 1\% to 99\%) \emph{Default:} 50\%. {\footnotesize\texttt{zxtune.core.aym.duty\_cycle}}
\item[3-channel mixer]  {\footnotesize\texttt{zxtune.sound.mixer.3}}
\end{description}
\subsection*{Z80 CPU (AY files with player code)}
\begin{description}
\item[Clockrate]  \emph{Choices:} ZX Spectrum 48K, Pentagon (3.5 MHz), ZX Spectrum 128K (3.5469 MHz), Amstrad CPC (4 MHz), Turbo (7 MHz). \emph{Default:} ZX Spectrum 48K, Pentagon (3.5 MHz). {\footnotesize\texttt{zxtune.core.z80.clockrate}}
\item[Interrupt duration] How long interrupt request is active each frame (range: 1 T to 100 T) \emph{Default:} 24 T. {\footnotesize\texttt{zxtune.core.z80.int\_ticks}}
\end{description}
\subsection*{YM2203 (TurboFM)}
\begin{description}
\item[Clockrate]  \emph{Choices:} 1.75 MHz, TurboFM (3.5 MHz), NTSC (3.58 MHz), 4 MHz. \emph{Default:} TurboFM (3.5 MHz). {\footnotesize\texttt{zxtune.core.fm.clockrate}}
\end{description}
\subsection*{YM2612/YM3438 (OPN2)}
\begin{description}
\item[Emulation core]  \emph{Choices:} Nuked (accurate, high CPU usage), ymfm (MAME), GPGX (fast). \emph{Default:} Nuked (accurate, high CPU usage). {\footnotesize\texttt{zxtune.core.fm.opn2.core}}
\item[Platform]  \emph{Choices:} Mega Drive model 1 (VA0-VA6, YM2612), Mega Drive model 1 VA7 (YM3438), Mega Drive model 2 (YM3438), Genesis 3 (YM3438), YM3438 without filter, YMF276 without filter, Custom. \emph{Default:} Mega Drive model 1 (VA0-VA6, YM2612). {\footnotesize\texttt{zxtune.core.fm.opn2.preset}}
\item[Chip]  \emph{Choices:} YM2612, YM3438, YMF276. \emph{Default:} YM2612. {\footnotesize\texttt{zxtune.core.fm.opn2.chip}}
\item[DAC ladder effect] Crossover distortion of discrete YM2612, audible on quiet sounds (Nuked and ymfm cores) \emph{Default:} on. {\footnotesize\texttt{zxtune.core.fm.opn2.ladder}}
\item[Low-pass filter] Cutoff of the board filter, 0 is off (range: 0 Hz to 24000 Hz) \emph{Default:} 3390 Hz. {\footnotesize\texttt{zxtune.core.fm.opn2.lowpass}}
\item[Low-pass slope]  \emph{Choices:} 6 dB/oct, 12 dB/oct. \emph{Default:} 6 dB/oct. {\footnotesize\texttt{zxtune.core.fm.opn2.lowpass\_order}}
\item[DC blocking] Cutoff of output capacitors, 0 keeps DC bias of the chip (range: 0 Hz to 200 Hz) \emph{Default:} 5 Hz. {\footnotesize\texttt{zxtune.core.fm.opn2.highpass}}
\item[Stereo width] (range: 0\% to 100\%) \emph{Default:} 100\%. {\footnotesize\texttt{zxtune.core.fm.opn2.stereo}}
\item[Output response]  {\footnotesize\texttt{ui.fm.opn2.response}}
\end{description}
\subsection*{YM2151 (OPM)}
\begin{description}
\item[Emulation core]  \emph{Choices:} Nuked (accurate, very high CPU usage), ymfm (MAME), MAME (legacy, fast). \emph{Default:} ymfm (MAME). {\footnotesize\texttt{zxtune.core.fm.opm.core}}
\item[Platform]  \emph{Choices:} YM3012 (stereo DAC: X68000, arcade), YM3014 (mono DAC), Custom. \emph{Default:} YM3012 (stereo DAC: X68000, arcade). {\footnotesize\texttt{zxtune.core.fm.opm.preset}}
\item[Low-pass filter] Cutoff of the board filter, 0 is off (range: 0 Hz to 24000 Hz) \emph{Default:} 0 Hz. {\footnotesize\texttt{zxtune.core.fm.opm.lowpass}}
\item[Low-pass slope]  \emph{Choices:} 6 dB/oct, 12 dB/oct. \emph{Default:} 12 dB/oct. {\footnotesize\texttt{zxtune.core.fm.opm.lowpass\_order}}
\item[DC blocking] Cutoff of output capacitors, 0 keeps DC bias of the chip (range: 0 Hz to 200 Hz) \emph{Default:} 0 Hz. {\footnotesize\texttt{zxtune.core.fm.opm.highpass}}
\item[Stereo width] (range: 0\% to 100\%) \emph{Default:} 100\%. {\footnotesize\texttt{zxtune.core.fm.opm.stereo}}
\item[Output response]  {\footnotesize\texttt{ui.fm.opm.response}}
\end{description}
\subsection*{YM2203/YM2608/YM2610 (OPN)}
\begin{description}
\item[Emulation core]  \emph{Choices:} ymfm (MAME), MAME (legacy). \emph{Default:} ymfm (MAME). {\footnotesize\texttt{zxtune.core.fm.opn.core}}
\item[Platform]  \emph{Choices:} PC-88/PC-98, Custom. \emph{Default:} PC-88/PC-98. {\footnotesize\texttt{zxtune.core.fm.opn.preset}}
\item[Low-pass filter] Cutoff of the board filter, 0 is off (range: 0 Hz to 24000 Hz) \emph{Default:} 0 Hz. {\footnotesize\texttt{zxtune.core.fm.opn.lowpass}}
\item[Low-pass slope]  \emph{Choices:} 6 dB/oct, 12 dB/oct. \emph{Default:} 12 dB/oct. {\footnotesize\texttt{zxtune.core.fm.opn.lowpass\_order}}
\item[DC blocking] Cutoff of output capacitors, 0 keeps DC bias of the chip (range: 0 Hz to 200 Hz) \emph{Default:} 0 Hz. {\footnotesize\texttt{zxtune.core.fm.opn.highpass}}
\item[Stereo width] (range: 0\% to 100\%) \emph{Default:} 100\%. {\footnotesize\texttt{zxtune.core.fm.opn.stereo}}
\item[Output response]  {\footnotesize\texttt{ui.fm.opn.response}}
\end{description}
\subsection*{YM3812/YMF262 (OPL)}
\begin{description}
\item[Emulation core]  \emph{Choices:} Nuked (accurate), AdLibEmu (fast), MAME (fast). \emph{Default:} Nuked (accurate). {\footnotesize\texttt{zxtune.core.fm.opl.core}}
\item[Platform]  \emph{Choices:} AdLib (mono), Sound Blaster 1.x/2.0, Sound Blaster Pro 1/2, Sound Blaster 16, Custom. \emph{Default:} Sound Blaster 16. {\footnotesize\texttt{zxtune.core.fm.opl.preset}}
\item[Low-pass filter] Cutoff of the board filter, 0 is off (range: 0 Hz to 24000 Hz) \emph{Default:} 0 Hz. {\footnotesize\texttt{zxtune.core.fm.opl.lowpass}}
\item[Low-pass slope]  \emph{Choices:} 6 dB/oct, 12 dB/oct. \emph{Default:} 6 dB/oct. {\footnotesize\texttt{zxtune.core.fm.opl.lowpass\_order}}
\item[DC blocking] Cutoff of output capacitors, 0 keeps DC bias of the chip (range: 0 Hz to 200 Hz) \emph{Default:} 0 Hz. {\footnotesize\texttt{zxtune.core.fm.opl.highpass}}
\item[Stereo width] (range: 0\% to 100\%) \emph{Default:} 100\%. {\footnotesize\texttt{zxtune.core.fm.opl.stereo}}
\item[Output response]  {\footnotesize\texttt{ui.fm.opl.response}}
\end{description}
\subsection*{YM2413 (OPLL)}
\begin{description}
\item[Emulation core]  \emph{Choices:} Nuked (accurate), EMU2413 (fast), MAME (fast). \emph{Default:} Nuked (accurate). {\footnotesize\texttt{zxtune.core.fm.opll.core}}
\item[Low-pass filter] Cutoff of the board filter, 0 is off (range: 0 Hz to 24000 Hz) \emph{Default:} 0 Hz. {\footnotesize\texttt{zxtune.core.fm.opll.lowpass}}
\item[Low-pass slope]  \emph{Choices:} 6 dB/oct, 12 dB/oct. \emph{Default:} 12 dB/oct. {\footnotesize\texttt{zxtune.core.fm.opll.lowpass\_order}}
\item[DC blocking] Cutoff of output capacitors, 0 keeps DC bias of the chip (range: 0 Hz to 200 Hz) \emph{Default:} 0 Hz. {\footnotesize\texttt{zxtune.core.fm.opll.highpass}}
\item[Output response]  {\footnotesize\texttt{ui.fm.opll.response}}
\item[Interpolation]  \emph{Default:} off. {\footnotesize\texttt{zxtune.core.dac.interpolation}}
\item[Clockrate]  \emph{Choices:} SAM Coupé (8 MHz), Creative Music System (7.16 MHz). \emph{Default:} SAM Coupé (8 MHz). {\footnotesize\texttt{zxtune.core.saa.clockrate}}
\item[Interpolation]  \emph{Choices:} None, Performance, Quality (high CPU usage). \emph{Default:} Performance. {\footnotesize\texttt{zxtune.core.saa.interpolation}}
\end{description}
\subsection*{Main screen}
\begin{description}
\item[Show track information] Title, chips, stream format, chip emulation, original driver and engine above the controls. Also toggled by the button at the left of the controls \emph{Default:} off. {\footnotesize\texttt{ui.brief.expanded}}
\item[Show author] Display author under the track title \emph{Default:} off. {\footnotesize\texttt{ui.brief.show\_author}}
\end{description}
\subsection*{Oscilloscope}
\begin{description}
\item[Oscilloscope] Display, trigger and performance {\footnotesize\texttt{ui.oscilloscope}}
\end{description}
\subsection*{Roland MT-32 (MIDI)}
\begin{description}
\item[ROMs folder] Select folder with MT-32 or CM-32L control and PCM ROM dumps (full or chip halves, detected by contents) {\footnotesize\texttt{zxtune.core.mt32.roms\_path}}
\item[Preferred model]  \emph{Choices:} Automatic (newest MT-32 first), MT-32, CM-32L / LAPC-I. \emph{Default:} Automatic (newest MT-32 first). {\footnotesize\texttt{zxtune.core.mt32.model}}
\item[Oscilloscope voices]  \emph{Choices:} 9 parts (8 melodic + rhythm), 32 partials (LA synthesis). \emph{Default:} 9 parts (8 melodic + rhythm). {\footnotesize\texttt{zxtune.core.mt32.voices}}
\end{description}
\subsection*{SNES (SPC, SNSF)}
\begin{description}
\item[Emulation engine] ares: cycle accurate, plays stereo as the console. snes\_spc: lighter, without phase inverted surround. Applied from the next track \emph{Choices:} snes\_spc, ares. \emph{Default:} ares. {\footnotesize\texttt{zxtune.core.spc.engine}}
\end{description}
\subsection*{MOS6581/SID}
\begin{description}
\item[Emulation engine]  \emph{Choices:} reSID, reSIDfp. \emph{Default:} reSIDfp. {\footnotesize\texttt{zxtune.core.sid.engine}}
\item[Default chip model] Used if tune does not specify it \emph{Choices:} MOS6581, MOS8580. \emph{Default:} MOS6581. {\footnotesize\texttt{zxtune.core.sid.model}}
\item[Force chip model] Ignore value specified in tune \emph{Default:} off. {\footnotesize\texttt{zxtune.core.sid.model\_force}}
\item[Default C64 model] (used if tune does not specify it) \emph{Choices:} PAL, NTSC, NTSC (old), Drean (PAL-N), PAL-M. \emph{Default:} PAL. {\footnotesize\texttt{zxtune.core.sid.clock}}
\item[Force C64 model] Ignore value specified in tune \emph{Default:} off. {\footnotesize\texttt{zxtune.core.sid.clock\_force}}
\item[Interpolation]  \emph{Choices:} None, Performance, Quality (high CPU usage). \emph{Default:} None. {\footnotesize\texttt{zxtune.core.sid.interpolation}}
\item[Filter emulation]  \emph{Default:} on. {\footnotesize\texttt{zxtune.core.sid.filter}}
\item[MOS6581 filter preset]  \emph{Choices:} Custom, JSIDPlay2 reSID: Lightest, JSIDPlay2 reSID: Lighter, JSIDPlay2 reSID: Light, JSIDPlay2 reSID: Average, JSIDPlay2 reSID: Dark, JSIDPlay2 reSID: Darker, JSIDPlay2 reSID: Darkest, VICE default (reSID bias +500 mV), ReSID 6581 (resid-fp fit), Alankila 6581R4AR 3789, Alankila 6581R3 3984 \#1, Alankila 6581R3 3984 \#2, Lord Nightmare 6581R3 4285, Lord Nightmare 6581R3 4485, Lord Nightmare 6581R4 1986S, ZrX 6581R3 0384, ZrX 6581R3 1984, ZrX 6581R3 3684, ZrX 6581R3 3985, ZrX 6581R4AR 2286, Trurl 6581R3 0784, Trurl 6581R3 0486S, Trurl 6581R3 3384, Trurl 6581R3 4885, Trurl 6581R4AR 3789, Trurl 6581R4AR 4486, Nata 6581R3 2083, Grue 6581R4AR 3488, KruLLo, Enigma 6581R3 4885, Enigma 6581R3 1585. \emph{Default:} Custom. {\footnotesize\texttt{zxtune.core.sid.filter\_6581\_preset}}
\item[MOS8580 filter preset]  \emph{Choices:} Custom, JSIDPlay2 reSID: Light, JSIDPlay2 reSID: Average, JSIDPlay2 reSID: Dark, Trurl 8580R5 1489, Trurl 8580R5 3691. \emph{Default:} Custom. {\footnotesize\texttt{zxtune.core.sid.filter\_8580\_preset}}
\item[Filter cutoff curve]  {\footnotesize\texttt{ui.sid.filter\_curve}}
\item[MOS8580 digi boost] MOS8580 only: make samples played via volume register audible \emph{Default:} off. {\footnotesize\texttt{zxtune.core.sid.digiboost}}
\item[reSID: MOS6581 filter bias] reSID engine, MOS6581 custom filter only: shifts filter cutoff frequency mapping (range: -1000 mV to 1000 mV) \emph{Default:} 0 mV. {\footnotesize\texttt{zxtune.core.sid.filter\_bias}}
\item[reSIDfp: MOS6581 filter curve] reSIDfp engine, custom filter only: 0 - light, 100 - dark \emph{Default:} 50. {\footnotesize\texttt{zxtune.core.sid.filter\_6581\_curve}}
\item[reSIDfp: MOS6581 filter range] reSIDfp engine, custom filter only: 0 - dark, 100 - light \emph{Default:} 50. {\footnotesize\texttt{zxtune.core.sid.filter\_6581\_range}}
\item[reSIDfp: MOS8580 filter curve] reSIDfp engine, custom filter only: 0 - light, 100 - dark \emph{Default:} 50. {\footnotesize\texttt{zxtune.core.sid.filter\_8580\_curve}}
\item[reSIDfp: combined waveforms]  \emph{Choices:} Average, Weak, Strong. \emph{Default:} Average. {\footnotesize\texttt{zxtune.core.sid.combined\_waveforms}}
\item[6581 pulse+saw (reSIDfp)]  \emph{Choices:} Normal (most 6581 chips), Strong (chips with loud pulse+saw). \emph{Default:} Normal (most 6581 chips). {\footnotesize\texttt{zxtune.core.sid.strong\_pulse\_saw}}
\item[2SID/3SID output]  \emph{Choices:} Stereo (chips panned), Mono. \emph{Default:} Stereo (chips panned). {\footnotesize\texttt{zxtune.core.sid.multi\_sid\_mono}}
\item[Samples lane (\$D418)] Oscilloscope: extra lane with master volume samples (4-bit digis), the three voices lanes show oscillators only. Applied from the next track \emph{Default:} off. {\footnotesize\texttt{zxtune.core.sid.digi\_lane}}
\end{description}
\subsection*{Android Auto}
\begin{description}
\item[Android Auto] Car display options {\footnotesize\texttt{ui.car}}
\end{description}
\subsection*{Playlist backup}
\begin{description}
\item[Playlist backup] Export or import the playlist, export a diagnostic log {\footnotesize\texttt{rescue}}
\end{description}
\section{Oscilloscope}
\subsection*{Display}
\begin{description}
\item[Waveforms window] Displayed duration of each voice waveform (range: 5 ms to 100 ms) \emph{Default:} 40 ms. {\footnotesize\texttt{ui.oscilloscope.window\_ms}}
\item[Visual amplification]  \emph{Choices:} Automatic (per voice), $\times$1, $\times$1.5, $\times$2, $\times$3, $\times$4, $\times$6, $\times$8. \emph{Default:} Automatic (per voice). {\footnotesize\texttt{ui.oscilloscope.gain}}
\item[Max periods shown (zoom on high notes)]  \emph{Choices:} Off (whole window), 4 periods, 8 periods, 16 periods, 32 periods. \emph{Default:} 8 periods. {\footnotesize\texttt{zxtune.sound.scope.max\_periods}}
\item[SID gauges window] Displayed duration of oscillators and master volume gauges (range: 2 ms to 40 ms) \emph{Default:} 10 ms. {\footnotesize\texttt{ui.oscilloscope.gauge\_window\_ms}}
\item[Oscilloscope height] Relative to 16:9 of width, the dashboard below takes the rest of screen (range: 60\% to 160\%) \emph{Default:} 100\%. {\footnotesize\texttt{ui.oscilloscope.height}}
\item[Filter response height] Part of oscilloscope taken by SID filter response above voices waveforms (range: 20\% to 80\%) \emph{Default:} 55\%. {\footnotesize\texttt{ui.oscilloscope.filter\_share}}
\item[AdLib rhythm lanes] OPL/OPL2: lanes of drums of rhythm mode (most of tunes don't use it), else 9 voices in 3$\times$3. Applied from the next track \emph{Default:} off. {\footnotesize\texttt{ui.oscilloscope.opl\_rhythm}}
\item[Waveforms under filter response] SID filter view shows response curve above voices waveforms, else response only \emph{Default:} on. {\footnotesize\texttt{ui.oscilloscope.filter\_waves}}
\item[Show technical statistics] Scope, sound output, rendering and GPU lines below the dashboard \emph{Default:} on. {\footnotesize\texttt{ui.oscilloscope.show\_stats}}
\item[Frame rate limit]  \emph{Choices:} Display refresh rate, 120 fps, 90 fps, 60 fps, 30 fps. \emph{Default:} Display refresh rate. {\footnotesize\texttt{ui.oscilloscope.max\_fps}}
\item[Graphics backend]  \emph{Choices:} Automatic (Vulkan if available), Vulkan, OpenGL ES. \emph{Default:} Automatic (Vulkan if available). {\footnotesize\texttt{ui.oscilloscope.backend}}
\item[Custom Vulkan driver] Select driver package (zip with meta.json or .so library, e.g. Mesa Turnip for Adreno GPUs) {\footnotesize\texttt{ui.oscilloscope.vulkan\_driver}}
\item[Use custom Vulkan driver] System driver is used if custom one fails to load \emph{Default:} off. {\footnotesize\texttt{ui.oscilloscope.vulkan\_driver.enabled}}
\end{description}
\subsection*{Trigger}
\begin{description}
\item[Preset]  \emph{Choices:} Default (corrscope template, 40 ms), Evolving notes (60 ms), Chiptune square/pulse, Bass, low notes, NES triangle, Power saving, Custom. \emph{Default:} Default (corrscope template, 40 ms). {\footnotesize\texttt{ui.oscilloscope.trigger\_preset}}
\item[Trigger width] History buffer size, should exceed two periods of the lowest notes (range: 10 ms to 200 ms) \emph{Default:} 40 ms. {\footnotesize\texttt{zxtune.sound.scope.trigger.width}}
\item[Edge strength] Prefer steep rising edges: new notes are caught better, but sustained ones jump more (range: 0\% to 300\%) \emph{Default:} 100\%. {\footnotesize\texttt{zxtune.sound.scope.trigger.edge\_strength}}
\item[Slope width] (range: 5\% of period to 100\% of period) \emph{Default:} 25\% of period. {\footnotesize\texttt{zxtune.sound.scope.trigger.slope\_width}}
\item[Buffer strength] Prefer similarity with previous frames (range: 0\% to 300\%) \emph{Default:} 100\%. {\footnotesize\texttt{zxtune.sound.scope.trigger.buffer\_strength}}
\item[Buffer responsiveness] (range: 0\% to 100\%) \emph{Default:} 50\%. {\footnotesize\texttt{zxtune.sound.scope.trigger.responsiveness}}
\item[Buffer falloff] (range: 10\% of period to 200\% of period) \emph{Default:} 50\% of period. {\footnotesize\texttt{zxtune.sound.scope.trigger.buffer\_falloff}}
\item[Reset below match] Forget history on new notes when wave doesn't match it (range: 0\% to 100\%) \emph{Default:} 30\%. {\footnotesize\texttt{zxtune.sound.scope.trigger.reset\_below}}
\item[DC removal rate] (range: 0\% to 100\%) \emph{Default:} 0\%. {\footnotesize\texttt{zxtune.sound.scope.trigger.mean\_responsiveness}}
\item[Sign amplification] Steep zero crossings for waves without clear edges (e.g. triangle) (range: 0\% to 300\%) \emph{Default:} 0\%. {\footnotesize\texttt{zxtune.sound.scope.trigger.sign\_strength}}
\item[Search range] (range: 10\% of width to 100\% of width) \emph{Default:} 50\% of width. {\footnotesize\texttt{zxtune.sound.scope.trigger.diameter}}
\item[Falling edges] E.g. for sampled trumpets \emph{Default:} off. {\footnotesize\texttt{zxtune.sound.scope.trigger.falling}}
\item[Pitch tracking] Rescale history on pitch changes \emph{Default:} on. {\footnotesize\texttt{zxtune.sound.scope.trigger.pitch\_tracking}}
\item[Zero crossing post trigger, samples] Slide to the nearest rising zero crossing, 0 to disable \emph{Default:} 0. {\footnotesize\texttt{zxtune.sound.scope.trigger.post\_radius}}
\item[Full searches]  \emph{Choices:} Every frame, 60 per second, 30 per second. \emph{Default:} Every frame. {\footnotesize\texttt{zxtune.sound.scope.trigger.search\_rate}}
\item[Trigger subsampling]  \emph{Choices:} Automatic, None (full rate), 2, 4. \emph{Default:} Automatic. {\footnotesize\texttt{zxtune.sound.scope.trigger.subsampling}}
\end{description}
\section{Mixer}
\subsection*{Channel A}
\begin{description}
\item[Left] (range: 0\% to 100\%) \emph{Default:} 100\%. {\footnotesize\texttt{zxtune.sound.mixer.3.0\_0}}
\item[Right] (range: 0\% to 100\%) \emph{Default:} 0\%. {\footnotesize\texttt{zxtune.sound.mixer.3.0\_1}}
\end{description}
\subsection*{Channel B}
\begin{description}
\item[Left] (range: 0\% to 100\%) \emph{Default:} 60\%. {\footnotesize\texttt{zxtune.sound.mixer.3.1\_0}}
\item[Right] (range: 0\% to 100\%) \emph{Default:} 60\%. {\footnotesize\texttt{zxtune.sound.mixer.3.1\_1}}
\end{description}
\subsection*{Channel C}
\begin{description}
\item[Left] (range: 0\% to 100\%) \emph{Default:} 0\%. {\footnotesize\texttt{zxtune.sound.mixer.3.2\_0}}
\item[Right] (range: 0\% to 100\%) \emph{Default:} 100\%. {\footnotesize\texttt{zxtune.sound.mixer.3.2\_1}}
\end{description}
\section{Android Auto}
\subsection*{Android Auto}
\begin{description}
\item[Show oscilloscope / dashboard in car] Replaces the cover art on the car screen with live visualization while playing. Read the legal notice before enabling. \emph{Default:} off. {\footnotesize\texttt{ui.car.visualization}}
\item[Car visualization content]  \emph{Choices:} Oscilloscope, Chips dashboard, Both. \emph{Default:} Oscilloscope. {\footnotesize\texttt{ui.car.visualization.mode}}
\item[Car visualization frame rate]  \emph{Choices:} 1 fps, 2 fps, 5 fps, 10 fps, 15 fps. \emph{Default:} 5 fps. {\footnotesize\texttt{ui.car.visualization.fps}}
\item[Legal notice] Rules about screens visible to the driver in your country {\footnotesize\texttt{ui.car.legal}}
\end{description}


\chapter{Technical statistics}
Below the dashboard, technical lines report on rendering:
\begin{description}
  \item[scope] lanes shown, trigger cost, synchronisation;
  \item[output] output rate and audio buffer size;
  \item[render] emulation speed relative to real time (above 100~\%, the device keeps up);
  \item[vulkan / gles] frames per second and drawing time;
  \item[vk] GPU and driver.
\end{description}

\chapter{Troubleshooting}
\begin{description}
  \item[Crackles or dropouts] is the ``render'' line below 100~\%? Choose a lighter engine
        (chapter~\ref{chap:emu-settings}), lower the oscilloscope frame rate or switch to Vulkan if available.
  \item[Stuttering oscilloscope] enable Vulkan, check that no battery saver throttles the application.
  \item[File not recognised] check it is in the formats table; xSF rips (PSF, SSF, GSF\dots) often need their
        \texttt{\_lib} files in the same folder.
  \item[Track too short or too long] set the default duration or the loops limit; for SID, the duration comes from the
        HVSC database.
  \item[Wrong notes on AY] set the clock to the original machine's.
  \item[No sound on MT-32] check that the module ROM images are present.
  \item[Catalogue unavailable] the third-party service may be offline; files already played stay cached.
  \item[Empty playlist after update] former XSPF playlists remain in the browser, from where they can be added to a
        playlist.
\end{description}

\section{Reporting a problem}
Give the version, the device, the file and, if a setting is involved, its internal key (chapter~\ref{chap:all-settings}).
The statistics lines below the dashboard are useful too.

\chapter{Known limitations}
We would rather tell you what does not work yet than let you find out.
\begin{itemize}
  \item Android Auto shows at most 32 lanes: beyond that, the last ones are not shown.
  \item Saturn SSF and PlayStation PSF files do not have a detailed per-voice view yet; their oscilloscope shows the
        overall output.
  \item The PC Engine view only serves VGM recordings; HES files use the generic view.
  \item Settings marked ``applied to the next track'' do not change the current track.
  \item The interface is in English only. Existing translations are kept in the sources but no longer shipped, since
        we cannot keep them up to date.
  \item The MT-32 needs the original module's memory images, which we do not supply.
  \item For now we only develop the Android version; the desktop version follows with a delay.
\end{itemize}


\part{Development}
\chapter{Architecture}\label{chap:architecture}

\section{Overview}
\soft has two layers:
\begin{description}
  \item[Native core (C++)] format parsing, chip emulation, mixing, per-voice outputs and chip states. It is built
        with the ZXTune \texttt{make} system for four Android architectures into \texttt{libzxtune.so}.
  \item[Application (Kotlin, Java)] playback service, interface, playlists, browsing, oscilloscope rendering (Vulkan
        or OpenGL ES) and dashboards.
\end{description}

\section{Audio path}
\begin{enumerate}
  \item A \emph{Module::Holder} is created for the file; it provides a \emph{Renderer} per playback.
  \item The renderer calls the emulation engine by blocks (about 100~ms), resamples to the output rate (48~kHz) and
        applies gain, fades and loop limits.
  \item The Android service writes the result to the audio output. Emulation state snapshots, taken regularly, make
        backward seeking fast without replaying from the start.
\end{enumerate}

\section{Visualisation data}
Every engine exposes the same three streams through the \emph{VoicesScope} interface:
\begin{description}
  \item[Layout] list of voice groups (one chip per group) and their names, sent on every change.
  \item[Per-voice outputs] samples of each voice at the output rate, resampled like the sound.
  \item[Voice states] snapshot of each voice about every millisecond: frequency, level, flags, chip kind and 64 bytes
        specific to the chip (registers, envelopes, wavetable\dots).
\end{description}
For libvgm chips these streams come from an extension added to the emulation cores: a per-sample callback (per-voice
outputs) and a state function. reSID and reSIDfp, ares and Ootake received equivalent extensions.

\section{Rendering}
\begin{itemize}
  \item A worker thread prepares each frame: for each voice, the time window is aligned on the fundamental period of
        the note (trigger), then reduced to two points per pixel column (minimum and maximum), so no peak is lost.
  \item Traces are drawn by Vulkan when available, otherwise OpenGL ES, with thin antialiased lines.
  \item The dashboard is drawn in two layers: a static layer (frames, captions), redrawn only when the layout
        changes, and a dynamic layer (values, curves), redrawn every frame.
\end{itemize}

\section{Tests}
\begin{itemize}
  \item More than 1~000 JVM tests (Robolectric) for the application: models, playlist storage, view computations, and
        render tests producing the screenshots of this manual.
  \item Native check tools: oscilloscope measurement (\texttt{scope\_bench}), per-voice outputs of VGM chips, SID
        parameters, seeking.
\end{itemize}

\chapter{Development heuristics}
Rules applied throughout development, and the reason for each.

\begin{description}
  \item[Fidelity first] the default engine is the most accurate available (reSIDfp, ares, Nuked, Ootake). A lighter
        engine remains selectable, for slow devices and comparison.
  \item[Exact values] every value shown is the chip state at the current frame. No smoothing, no averaging: an
        averaged value is a value that never existed.
  \item[Graphics when enough] when a curve says it all (filter response), no text repeats it. When text is needed,
        it sits next to what it describes.
  \item[One view per synthesis] each chip family has its own view, built around what makes its sound: wavetable for
        the PC Engine, operators and algorithms for FM, envelopes and echo for the Super Nintendo, filter and
        modulations for the SID.
  \item[Measure before claiming] every fix is checked by a measurement (compared per-voice outputs, render, test)
        before delivery; what could not be checked is reported as such.
  \item[Render before delivery] every visual change is checked on a render generated by the tests before being built
        into an APK.
  \item[Minimal patches to third-party code] changes to imported libraries are kept small, marked in the code
        (``zxtune extension'', ``standalone'') and follow the style of the modified file, so they can be merged back.
  \item[Neutral defaults] features that change the layout (SID samples lane, AdLib rhythm lanes) are off by default.
  \item[Licenses checked on import] all imported code is checked for compatibility with GPL~v3.
  \item[No regression] the full test suite runs before significant deliveries; the only tolerated failures are
        catalogue tests needing network access.
\end{description}

\chapter{Development cost estimate}\label{chap:chiffrage}

We wanted to know what this work would have cost had it been given to an experienced freelance developer, familiar
with emulation and Android. This estimate covers only what \soft adds to ZXTune between 2 October and
24 November 2026.

\textbf{The cost of the proprietary NebulaWeave drivers is not included and will not be disclosed in this document:
it is a trade secret.}

\section{Why we do not count lines}
A first estimate, based on the number of lines added to the repository, gave more than 900 days of work. That figure
does not hold: a good share of those lines is code imported as is (ymfm, mt32emu, ares), resources, tests generated
in series, or code longer than it needed to be. A developer who knows the subject writes less code for the same
result. So we estimated feature by feature, asking each time how many days a competent person would need to get the
same behaviour, tests included.

We also separate created code from adapted code. Adapting an existing emulation core, for instance adding to Ootake
the reading of each voice's wavetable, is mostly a matter of understanding the core; the code added fits in a few
hundred lines. Integrating an imported library only takes checking its licence, building it for Android and wiring it
in.

\section{Breakdown}
\begin{longtable}{@{}p{0.75\textwidth} r@{}}
\toprule \textbf{Feature} & \textbf{Days} \\ \midrule
Per-voice outputs and chip states in the engines (libvgm, SID, SNES, PC Engine, PCM) & 25 \\
Vulkan and OpenGL ES oscilloscope, triggering & 15 \\
SID dashboard and filter response & 15 \\
FM views and inspector & 12 \\
Super Nintendo view & 6 \\
PCM chip view (YMW258, SCSP) & 6 \\
PC Engine view & 4 \\
NES and MT-32 cards & 3 \\
ares port and SNSF integration & 6 \\
FM core choice, presets, output stage & 8 \\
SID options (samples lane, mono, filters) & 4 \\
Multiple playlists, favourites, queue, M3U & 8 \\
Interface redesign & 15 \\
Audio formats and miscellaneous fixes & 2 \\
Desktop version (Qt) & 12 \\
VGM extraction tools (Model 2 and 3) & 8 \\
Tests and control renders & 10 \\
Manual & 6 \\ \midrule
Subtotal & 165 \\
Design, research, acceptance (15~\%) & 25 \\ \midrule
\textbf{Total} & \textbf{190} \\
\bottomrule
\end{longtable}

\section{Cost}
\begin{longtable}{@{}l r r@{}}
\toprule \textbf{Assumption} & \textbf{Day rate} & \textbf{Cost} \\ \midrule
Low & 450~€ & 85\,500~€ \\
Median & 600~€ & 114\,000~€ \\
High & 750~€ & 142\,500~€ \\
\bottomrule
\end{longtable}

These amounts include neither the ZXTune base, developed over years by its author, nor the free libraries reused
unchanged, nor the proprietary drivers. They give an order of magnitude: adjust them if your day rate or pace differs
from ours.


\appendix
\chapter{Glossary}
\begin{description}
  \item[ADSR] attack, decay, sustain, release: the four phases of a note envelope.
  \item[Cent] one hundredth of a semitone.
  \item[DDA] HuC6280 mode where the processor writes samples directly.
  \item[DSP] digital signal processor (echo, reverb).
  \item[FM] frequency modulation synthesis.
  \item[LFSR] feedback shift register producing noise.
  \item[Operator] sine oscillator of an FM voice.
  \item[PAL-N] Argentine video standard: clock close to NTSC, 50 frames per second.
  \item[Wavetable] stored waveform played in a loop.
  \item[Tracker] pattern-based composition software, and the format of the music it produces.
\end{description}

\chapter{Version history}
This appendix summarises the main steps of \soft development, most recent first. Numbers are those of the
development builds delivered for testing.

\section*{Version 40}
\begin{itemize}
  \item Interface entirely in English.
  \item SID: the Mod column names the source voice only (``Sync V1'', ``Sync+Ring V1'').
\end{itemize}
\section*{Version 39}
\begin{itemize}
  \item PC Engine view: waveform shape named (sine, triangle, rising or falling saw, square, pulse), table number and
        range.
\end{itemize}
\section*{Version 38}
\begin{itemize}
  \item PC Engine view: level bar and panning bars.
  \item SID: the samples lane sits under the voices of its chip; the three voices no longer contain samples played
        through the volume.
  \item AdLib OPL/OPL2: nine voices in 3~$\times$~3, drum lanes hidden by default.
\end{itemize}
\section*{Version 37}
\begin{itemize}
  \item PC Engine: per-voice output and HuC6280 state in the Ootake core; dedicated view with each voice's wavetable
        and all its registers.
\end{itemize}
\section*{Version 36}
\begin{itemize}
  \item Fullscreen portrait: top cutout and bottom corners avoided, dashboard correctly placed.
  \item SID: the Mod column names the source voice of sync and ring.
\end{itemize}
\section*{Version 35}
\begin{itemize}
  \item Oscilloscope height and filter response height settings.
  \item SID: optional samples lane (\$D418).
  \item Live panel: SNES page.
\end{itemize}
\section*{Versions 33 and 34}
\begin{itemize}
  \item Multiple playlists, favourites, adding from the player, the browser and playlists.
  \item Play next queue, M3U import and export, $\equiv$+ button in the player.
  \item Multi-song files played directly from the browser.
  \item SID: optional mono output for 2SID and 3SID tunes.
\end{itemize}
\section*{Versions 30 to 32}
\begin{itemize}
  \item Super Nintendo: ares (cycle-accurate) emulation by default, snes\_spc as an option, for SPC and SNSF.
\end{itemize}
\section*{Earlier versions}
\begin{itemize}
  \item SID filter response view in the oscilloscope, without value texts.
  \item SCSP effects DSP enabled: reverb of Sega arcade boards and the Saturn.
  \item Dashboard values in real time, frame by frame.
  \item Technical track information (codec, rate, bit depth, bitrate), compact bottom area.
  \item Selectable FM cores per family, platform presets and analogue output stage.
  \item YMW258 and SCSP PCM chip view.
  \item Extended WAV (PCM and floating point).
  \item LHUNA identity: navigation bar, mini player, settings as cards, About screen.
  \item Vulkan and OpenGL ES oscilloscope with thin traces, window fitted to each voice period.
  \item FM inspector (OPN, OPM, OPL) and S-DSP view.
\end{itemize}

\chapter{Credits and component licenses}\label{chap:credits}
\soft relies on the work of many free software projects. The table below is generated from the license files present
in the sources; the full text of each license comes with the source code and is shown in the About screen of the
application.

Application base: \texttt{zxtune}, under LGPL~v3.
% generated by tools/credits.py, do not edit
\begin{longtable}{@{}l l@{}}
\toprule \textbf{Composant (dossier)} & \textbf{Licence} \\ \midrule \endhead
\texttt{FLAC} & BSD (Xiph) \\
\texttt{adrenotools} & BSD 2-clause \\
\texttt{asap} & GPL v2 \\
\texttt{atrac9} & MIT \\
\texttt{curl} & curl (MIT-like) \\
\texttt{ffmpeg} & LGPL v2.1 \\
\texttt{gme/gme} & LGPL v2.1 \\
\texttt{he} & GPL v3 \\
\texttt{ht} & GPL v3 \\
\texttt{ht/Core} & GPL v3 \\
\texttt{hvl} & BSD 3-clause \\
\texttt{lame} & LGPL v2 \\
\texttt{lhasa} & ISC \\
\texttt{mgba} & MPL 2.0 \\
\texttt{minimp3} & CC0 \\
\texttt{mpg123} & LGPL v2.1 \\
\texttt{ogg} & BSD \\
\texttt{openmpt} & BSD \\
\texttt{opus} & BSD \\
\texttt{sidplayfp} & GPL v2 \\
\texttt{snesspc} & LGPL v2.1 \\
\texttt{snesspc/ares} & ISC \\
\texttt{sseqplayer} & WTFPL \\
\texttt{unrar} & UnRAR (freeware, source available) \\
\texttt{v2m} & Artistic 2.0 \\
\texttt{vgmstream} & ISC-like (see COPYING) \\
\texttt{vio2sf/desmume} & GPL v2 \\
\texttt{vorbis} & BSD \\
\texttt{xmp} & MIT \\
\texttt{z80ex} & GPL v2 \\
\bottomrule \end{longtable}

The libvgm library (\texttt{vgm} folder) and the cores it contains (Nuked, MAME, Ootake, EMU2413\dots) have no single
license file in the sources: the license of each core must be checked individually before publishing the APK under
GPL~v3.

\capturedouble{about-credits-en}{about-licenses}{Credits in the application}{licenses}
\section*{Main authors}
\begin{description}
  \item[ZXTune] Vitamin, base of the application.
  \item[libsidplayfp] Simon White, Antti Lankila, Leandro Nini and contributors; reSID by Dag Lem.
  \item[ares] Near and the ares team.
  \item[snes\_spc, Game Music Emu] Shay Green (blargg) and contributors.
  \item[libvgm] Valley Bell; Nuked cores by Alexey Khokholov, ymfm by Aaron Giles, Ootake by Kitao Nakamura.
  \item[vgmstream] Adam Gashlin, bnnm and contributors.
  \item[libopenmpt] the OpenMPT team; \textbf{libxmp} Claudio Matsuoka and contributors.
  \item[Munt] the Munt team.
  \item[ASAP] Piotr Fusik and contributors.
  \item[Highly Experimental / Theoretical] Neill Corlett (PlayStation and Saturn cores).
\end{description}
Project websites are listed in the bibliography.

\chapter{Bibliography}

\section*{Software projects}
\begin{itemize}
  \item ZXTune, by Vitamin : \url{https://github.com/vitamin-caig/zxtune}
  \item libsidplayfp, reSID et reSIDfp : \url{https://github.com/libsidplayfp/libsidplayfp}
  \item ares : \url{https://ares-emu.net}
  \item libvgm, by Valley Bell : \url{https://github.com/ValleyBell/libvgm}
  \item Nuked cores, by nukeykt : \url{https://github.com/nukeykt}
  \item ymfm, by Aaron Giles : \url{https://github.com/aaronsgiles/ymfm}
  \item vgmstream : \url{https://github.com/vgmstream/vgmstream}
  \item libopenmpt : \url{https://lib.openmpt.org}
  \item Munt (MT-32 emulator) : \url{https://github.com/munt/munt}
  \item ASAP (POKEY) : \url{https://asap.sourceforge.net}
\end{itemize}

\section*{Collections}
\begin{itemize}
  \item High Voltage SID Collection : \url{https://www.hvsc.c64.org}
  \item The Mod Archive : \url{https://modarchive.org}
  \item VGMRips : \url{https://vgmrips.net}
  \item ZXArt : \url{https://zxart.ee}
\end{itemize}

\section*{Chip documentation}
\begin{itemize}
  \item Manufacturer datasheets (General Instrument AY-3-8910, MOS 6581, Yamaha YM2612, YM3812, YMF262,
        YM2151, YMF292) and documentation of the emulation projects above.
  \item Community hardware documentation : NESdev Wiki (NES), fullsnes (Super Nintendo), Pan Docs
        (Game Boy), PC Engine scene documentation.
\end{itemize}

\section*{References for the settings}
\begin{itemize}
  \item JSIDPlay2 (SID filter presets) : \url{https://sourceforge.net/projects/jsidplay2/}
  \item VICE (reSID filter bias) : \url{https://vice-emu.sourceforge.io}
  \item sidplayfp, manual page \texttt{sidplayfp.ini} : \url{https://www.mankier.com/5/sidplayfp.ini}
  \item DOSBox Staging, output filter emulation : \url{https://www.dosbox-staging.org/0.83/manual/sound/output-filters/}
  \item BlastEm : \url{https://www.retrodev.com/blastem/}
  \item Genesis Plus GX, low-pass filter discussion : \url{https://forums.libretro.com/t/genesis-plus-gx-where-is-the-audio-lowpass-filter/12169}
  \item Console Mods, Genesis audio circuit : \url{https://consolemods.org/wiki/Genesis:Audio_Circuit_Mod}
  \item Mesen, SNES audio options : \url{https://www.mesen.ca/snes/docs/configuration/audio.html}
  \item Snes9x (cœur libretro), interpolation : \url{https://docs.libretro.com/library/snes9x/}
  \item corrscope (correlation trigger) : \url{https://github.com/corrscope/corrscope}
\end{itemize}


\end{document}
